% ─────────────────────────────────────────────────────────────────────────────
%  Monografia de TCC — Chatbot de Apoio Clínico para Manejo da ILTB
%  Autor      : Yago Novaes Neves
%  Orientador : <preencher>
%  LabPEC / DTI / UFES
%  Classe     : abnTeX2 (monografia, 12pt, ABNT NBR 14724)
%  Citações   : ABNT NBR 6023 (sistema autor-data, ALF)
% ─────────────────────────────────────────────────────────────────────────────

\documentclass[
  12pt,
  oneside,
  a4paper,
  english,
  brazil
]{abntex2}

% ─── Codificação e fonte ─────────────────────────────────────────────────────
\usepackage[T1]{fontenc}
\usepackage[utf8]{inputenc}
\usepackage{mathptmx}           % Times New Roman

% ─── Pacotes essenciais ───────────────────────────────────────────────────────
\usepackage{indentfirst}
\usepackage{color}
\usepackage[dvipsnames]{xcolor,colortbl}
\usepackage{graphicx}
\usepackage{ifthen}
\usepackage{array}
\usepackage{amssymb}
\usepackage{amsmath}
\usepackage{tabularx}
\usepackage{longtable}
\usepackage{booktabs}
\usepackage[singlelinecheck=false]{caption}
\usepackage{subcaption}
\usepackage{threeparttable}
\usepackage{enumitem}
\usepackage{xspace}
\usepackage{newfloat}
\DeclareFloatingEnvironment[
  name={Quadro},
  listname={Lista de Quadros},
  placement=htbp
]{quadro}
% ─── Citações ABNT ───────────────────────────────────────────────────────────
\usepackage[alf,
  abnt-full-initials=no,
  abnt-emphasize=bf,
  abnt-and-type=e,
  abnt-url-package=url,
  abnt-repeated-author-omit=yes,
  abnt-etal-cite=3,
  abnt-etal-list=3
]{abntex2cite}

\newcommand{\bibtextitlecommand}[2]{%
\ifthenelse{\equal{#1}{inproceedings}}{\normalfont{\textbf{#2} $[\ldots]$}}{}}

\usepackage{url6023}

% ─── Hyperlinks (abntex2 já carrega hyperref — não recarregar) ───────────────
\hypersetup{
  hidelinks,
  pdfsubject={Monografia de TCC — Chatbot RAG para Apoio ao Manejo da ILTB},
  pdfcreator={LaTeX with abnTeX2},
  pdfkeywords={RAG}{LLM}{ILTB}{tuberculose latente}{enfermagem}{LGPD},
  bookmarksdepth=4
}

% ─── Marcadores de revisão (implementação leve — sem TikZ/todonotes) ─────────
% \aluno{} aparece em verde no rascunho; para a versão final, comente as duas
% linhas abaixo e descomente as duas seguintes (que tornam os marcadores invisíveis).
\newcommand{\aluno}[1]{{\color{olive!80!black}\textbf{[\,}{\small\itshape PREENCHER: #1}\textbf{\,]}}}
\newcommand{\professor}[1]{{\color{red!70!black}\textbf{[\,}{\small\itshape PROF: #1}\textbf{\,]}}}
% Versão final (sem marcadores visíveis) — descomentar e comentar as linhas acima:
% \newcommand{\aluno}[1]{}
% \newcommand{\professor}[1]{}

% ─── Informações da monografia ──────────────────────────────────────────────
% ATENÇÃO: \titulo, \orientador, \instituicao e \preambulo são usados dentro
% de boxes na capa/folha de rosto. Não usar \aluno{} aqui — substitua os
% textos em MAIÚSCULAS pelos valores reais antes de compilar para entrega.
\titulo{Chatbot de Apoio Clínico para o Manejo da Infecção Latente por
  \textit{Mycobacterium tuberculosis} (ILTB): uma abordagem baseada em
  Recuperação com Geração Aumentada (RAG)}
\autor{Yago Novaes Neves}
\local{Vitória}
\data{2026}
\orientador{Prof. Dr. Wilian Hiroshi Hisatugu}
\coorientador{Prof. Dr. Renato Elias Nunes de Moraes}
\instituicao{%
  Universidade Federal do Espírito Santo -- UFES \\
  Centro Tecnológico -- Departamento de Engenharia de Produção (LabPEC/DTI)}
\tipotrabalho{Trabalho de Conclusão de Curso}
\preambulo{Trabalho de Conclusão de Curso apresentado ao Curso de Engenharia
  de Produção da Universidade Federal do Espírito Santo como requisito parcial
  para obtenção do grau de Bacharel em Engenharia de Produção.}

% ─────────────────────────────────────────────────────────────────────────────
\begin{document}
\selectlanguage{brazil}
\frenchspacing

% ═════════════════════════════════════════════════════════════════════════════
% ELEMENTOS PRÉ-TEXTUAIS
% ═════════════════════════════════════════════════════════════════════════════
\pretextual

% ─── Capa ────────────────────────────────────────────────────────────────────
\imprimircapa

% ─── Folha de rosto ──────────────────────────────────────────────────────────
\imprimirfolhaderosto

% ─── Folha de aprovação ──────────────────────────────────────────────────────
\begin{folhadeaprovacao}
  \begin{center}
    {\ABNTEXchapterfont\large\imprimirautor}

    \vspace*{\fill}\vspace*{\fill}

    {\ABNTEXchapterfont\bfseries\Large\imprimirtitulo}

    \vspace*{\fill}

    \imprimirpreambulo

    \vspace*{\fill}

    Vitória, \_\_\_\_ de \_\_\_\_\_\_\_\_\_\_ de 2026.

    \vspace*{\fill}

    \assinatura{\textbf{\imprimirorientador} \\ Orientador --- UFES}
    \assinatura{\textbf{\imprimircoorientador} \\ Coorientador --- UFES}
    \assinatura{\textbf{NOME DO MEMBRO EXAMINADOR} \\ INSTITUIÇÃO}

    \vspace*{\fill}
  \end{center}
\end{folhadeaprovacao}

% ─── Dedicatória (opcional) ──────────────────────────────────────────────────
% \begin{dedicatoria}
%   \vspace*{\fill}
%   \centering\noindent
%   \aluno{Texto de dedicatória.}
% \end{dedicatoria}

% ─── Agradecimentos (opcional) ───────────────────────────────────────────────
% \begin{agradecimentos}
%   \aluno{Texto de agradecimentos.}
% \end{agradecimentos}

% ─── Resumo em português ─────────────────────────────────────────────────────
\begin{resumo}

\aluno{Escrever resumo informativo (150--500 palavras): contextualização, justificativa,
objetivo, metodologia (DSRM), principais resultados (RAGAS gate final:
faithfulness 0,761, context\_precision 0,964), avaliação expert e conclusão.}

A Infecção Latente por \textit{Mycobacterium tuberculosis} (ILTB) representa um
desafio persistente para a saúde pública, com risco de progressão para tuberculose
ativa em grupos vulneráveis. Profissionais de enfermagem atuam na linha de frente do
rastreamento e tratamento da ILTB, mas frequentemente enfrentam lacunas no acesso a
informações clínicas atualizadas durante o atendimento. Este trabalho propõe,
desenvolve e avalia um chatbot de apoio clínico baseado na arquitetura de
Recuperação com Geração Aumentada (RAG), capaz de responder a perguntas sobre o
manejo da ILTB a partir de documentos oficiais do Ministério da Saúde e da
Organização Mundial da Saúde. Seguindo a \textit{Design Science Research
Methodology} (DSRM), o artefato foi construído com FastAPI, ChromaDB, embedding
multilíngue \texttt{BAAI/bge-m3} (1024D), busca híbrida BM25+denso com
\textit{Reciprocal Rank Fusion}, \textit{reranking} \textit{cross-encoder} mGTE
e o modelo de linguagem Llama~3.3~70B via API Groq. A avaliação automatizada
com o \textit{framework} RAGAS, conduzida sobre 38 questões clínicas, obteve no
gate final \textit{faithfulness} de 0{,}761 e \textit{context precision} de
0{,}964 (PASS sobre a meta de 0{,}75), salto de $+47{,}8\%$ em
\textit{faithfulness} frente ao \textit{baseline} pré-otimização.
\aluno{Completar com resultado da avaliação expert e conclusão principal.}

\vspace{\onelineskip}
\noindent\textbf{Palavras-chave:} tuberculose latente; ILTB; chatbot clínico;
recuperação com geração aumentada; RAG; grandes modelos de linguagem; enfermagem;
LGPD.
\end{resumo}

% ─── Abstract (em inglês) ────────────────────────────────────────────────────
\begin{resumo}[Abstract]
\begin{otherlanguage*}{english}

\aluno{Traduzir resumo para inglês após versão portuguesa finalizada.}

\vspace{\onelineskip}
\noindent\textbf{Keywords:} latent tuberculosis infection; LTBI; clinical chatbot;
retrieval-augmented generation; RAG; large language models; nursing; data privacy.
\end{otherlanguage*}
\end{resumo}

% ─── Lista de figuras ────────────────────────────────────────────────────────
\pdfbookmark[0]{\listfigurename}{lof}
\listoffigures*
\cleardoublepage

% ─── Lista de quadros ────────────────────────────────────────────────────────
\pdfbookmark[0]{Lista de Quadros}{loq}
\listof{quadro}{Lista de Quadros}
\cleardoublepage

% ─── Lista de tabelas ────────────────────────────────────────────────────────
\pdfbookmark[0]{\listtablename}{lot}
\listoftables*
\cleardoublepage

% ─── Lista de abreviaturas e siglas ──────────────────────────────────────────
\begin{siglas}
  \item[ABNT]   Associação Brasileira de Normas Técnicas
  \item[ANN]    Approximate Nearest Neighbor (busca por vizinhos aproximados)
  \item[API]    Application Programming Interface
  \item[BERT]   Bidirectional Encoder Representations from Transformers
  \item[CEP]    Comitê de Ética em Pesquisa
  \item[DSRM]   \textit{Design Science Research Methodology}
  \item[HNSW]   Hierarchical Navigable Small World
  \item[IGRA]   Interferon-Gamma Release Assay
  \item[ILTB]   Infecção Latente por \textit{Mycobacterium tuberculosis}
  \item[LLM]    \textit{Large Language Model} (grande modelo de linguagem)
  \item[LGPD]   Lei Geral de Proteção de Dados Pessoais
  \item[MS]     Ministério da Saúde
  \item[NLP]    \textit{Natural Language Processing} (processamento de linguagem natural)
  \item[OMS]    Organização Mundial da Saúde
  \item[PPD]    Derivado Proteico Purificado (Teste Tuberculínico)
  \item[PT]     Prova Tuberculínica
  \item[PVHIV]  Pessoas Vivendo com HIV
  \item[RAG]    \textit{Retrieval-Augmented Generation}
  \item[RAGAS]  \textit{Retrieval Augmented Generation Assessment}
  \item[REST]   \textit{Representational State Transfer}
  \item[SBERT]  Sentence-BERT
  \item[SUS]    \textit{System Usability Scale}
  \item[TB]     Tuberculose
  \item[UFES]   Universidade Federal do Espírito Santo
\end{siglas}

% ─── Sumário ─────────────────────────────────────────────────────────────────
\pdfbookmark[0]{\contentsname}{toc}
\tableofcontents*
\cleardoublepage

% ═════════════════════════════════════════════════════════════════════════════
% ELEMENTOS TEXTUAIS
% ═════════════════════════════════════════════════════════════════════════════
\textual

% ─────────────────────────────────────────────────────────────────────────────
\chapter{Introdução}
\label{cap:introducao}
% ─────────────────────────────────────────────────────────────────────────────

A tuberculose (TB) permanece como um dos maiores desafios de saúde pública no
mundo. Em 2022, a Organização Mundial da Saúde estimou 10,6 milhões de novos casos
de TB ativa, com o Brasil figurando entre os 30 países com maior carga da doença
\cite{who2023}. A forma latente da infecção --- denominada ILTB --- caracteriza-se
pela presença de \textit{Mycobacterium tuberculosis} no organismo sem manifestação
de doença ativa, mas com risco de progressão para TB ativa ao longo da vida,
sobretudo em indivíduos imunossuprimidos \cite{who2018}. O controle da ILTB é,
portanto, estratégico para a eliminação da TB como problema de saúde pública.

Nesse cenário, profissionais de enfermagem desempenham papel central no rastreamento,
na testagem e no acompanhamento de pacientes com ILTB. Contudo, os protocolos
clínicos são extensos, frequentemente atualizados e nem sempre de fácil consulta
durante o atendimento. Ferramentas de apoio à decisão clínica baseadas em
inteligência artificial (IA) surgem como alternativa promissora para reduzir essa
lacuna \cite{SHTI2023}.

Os grandes modelos de linguagem (\textit{Large Language Models} --- LLMs)
demonstraram capacidade notável de compreensão e geração de texto em domínios
especializados \cite{brown2020}. No entanto, modelos sem acesso a bases de
conhecimento atualizadas tendem a produzir respostas incorretas ou fabricadas ---
fenômeno conhecido como alucinação \cite{bang2023}. A arquitetura RAG
(\textit{Retrieval-Augmented Generation}) mitiga esse problema ao recuperar trechos
relevantes de documentos específicos antes de gerar a resposta \cite{lewis2021,
gao2024}.

\section{Problema e Motivação}

Apesar do avanço dos LLMs e da arquitetura RAG, três lacunas convergem para
caracterizar o problema enfrentado neste trabalho. A primeira é \textbf{linguística e
de domínio}: poucos trabalhos abordam aplicações de RAG em português do Brasil para
doenças infecciosas, e benchmarks como o PROPOR/ENAMED mostram que LLMs apresentam
desempenho inferior em consultas ao protocolo clínico brasileiro quando comparados
ao desempenho em corpora em inglês \cite{teachingLLMsBrazil2026,
adaptingLLMsPortuguese2024, llmsMedicalExam2026}.

A segunda é \textbf{regulatória e ética}: a Lei Geral de Proteção de Dados Pessoais
(LGPD) classifica dados de saúde como sensíveis e impõe restrições ao tratamento por
operadores terceirizados, o que problematiza arquiteturas RAG que dependem de APIs
externas para embeddings ou inferência \cite{lgpdSaude2023, lgpdEnfermagem2022,
sokPrivacyLLM2026, privacyRAGHealthcare2025}. Soluções voltadas à atenção primária do
SUS precisam reconciliar capacidade clínica com conformidade legal.

A terceira é \textbf{operacional}: a adoção de ferramentas de apoio à decisão
clínica por profissionais de enfermagem está condicionada a fatores de utilidade
percebida, alinhamento ético e prontidão organizacional --- aspectos modelados pelo
TAM-AIN (\textit{Technology Acceptance Model for AI in Nursing})
\cite{nursingTAM2024}. A ausência de instrumentação que cubra simultaneamente
acurácia técnica e essas dimensões de adoção é uma lacuna metodológica relevante
para o cenário brasileiro.

Este trabalho atende essas três lacunas ao desenvolver e avaliar um artefato RAG
para ILTB com corpus em português, pipeline de embeddings local (sem envio de
consultas a terceiros antes da busca) e protocolo de avaliação que combina métricas
RAGAS automáticas, rubrica clínica multidimensional e TAM-AIN.

\section{Objetivo Geral}

Desenvolver e avaliar um chatbot de apoio clínico ao manejo da ILTB baseado na
arquitetura RAG, com corpus em língua portuguesa, respeitando boas práticas de
privacidade de dados e avaliado com métricas objetivas de qualidade de resposta.

\section{Objetivos Específicos}

\begin{enumerate}[label=\textbf{OE\arabic*}]
  \item Construir e indexar um corpus de documentos clínicos sobre ILTB em
        português, provenientes de fontes oficiais;
  \item Implementar um pipeline RAG com embeddings multilíngues e LLM integrado
        via API;
  \item Avaliar o sistema com o \textit{framework} RAGAS em termos de
        \textit{faithfulness}, \textit{answer relevancy}, \textit{context precision}
        e \textit{context recall};
  \item Conduzir avaliação qualitativa com especialista em ILTB para verificar
        utilidade clínica e aderência aos protocolos vigentes.
\end{enumerate}

\section{Contribuições}

As contribuições deste trabalho são:

\begin{enumerate}[label=\textbf{C\arabic*}]
  \item \textbf{Artefato RAG aberto para ILTB em português:} pipeline RAG
        funcional, indexando documentos oficiais do MS e da OMS, com código-fonte
        disponível para reprodução e adaptação.
  \item \textbf{Dataset de avaliação clínica:} 38 questões clínicas em português
        sobre ILTB, organizadas em seis categorias (esquemas terapêuticos,
        monitoramento, populações especiais, diagnóstico, interações medicamentosas
        e efeitos adversos), com \textit{ground truths} extraídos literalmente dos
        protocolos indexados, utilizáveis como \textit{benchmark} por trabalhos
        subsequentes em RAG clínico em PT-BR.
  \item \textbf{Análise empírica do \textit{trade-off} LGPD vs. desempenho:}
        documentação das escolhas arquiteturais que viabilizam conformidade com a
        LGPD em uma POC acadêmica de saúde, e discussão das mitigações adicionais
        necessárias para implantação institucional.
  \item \textbf{Protocolo de avaliação multidimensional para RAG clínico com
        $N=1$:} adaptação metodológica que combina avaliação RAGAS automatizada
        com rubrica clínica multidimensional e TAM-AIN, oferecendo alternativa a
        instrumentos inadequados a amostra reduzida (e.g., Kappa de Cohen).
\end{enumerate}

\section{Estrutura da Monografia}

Esta monografia está organizada da seguinte forma. O
Capítulo~\ref{cap:referencial} apresenta o referencial teórico. O
Capítulo~\ref{cap:metodologia} descreve a metodologia DSRM adotada. O
Capítulo~\ref{cap:desenvolvimento} detalha o desenvolvimento do artefato. O
Capítulo~\ref{cap:avaliacao} apresenta a avaliação e os resultados. O
Capítulo~\ref{cap:discussao} discute as implicações dos resultados. O
Capítulo~\ref{cap:conclusao} conclui o trabalho com contribuições e trabalhos
futuros.

% ─────────────────────────────────────────────────────────────────────────────
\chapter{Referencial Teórico}
\label{cap:referencial}
% ─────────────────────────────────────────────────────────────────────────────

\section{Tuberculose e Infecção Latente (ILTB)}

A ILTB é definida pela OMS como um estado de resposta imune persistente à
estimulação por antígenos de \textit{Mycobacterium tuberculosis}, sem evidência de
TB ativa clinicamente manifesta \cite{who2018}. Estima-se que aproximadamente um
quarto da população mundial tenha sido infectada, com risco de progressão para TB
ativa de 5--10\% ao longo da vida em imunocompetentes \cite{who2023}. No Brasil, o
Ministério da Saúde define os critérios de rastreamento, diagnóstico e tratamento
por meio de seu Manual de Recomendações \cite{brasil2022}.

Determinados grupos apresentam risco aumentado de progressão para TB ativa e
constituem público prioritário para rastreamento e tratamento da ILTB. Entre eles
destacam-se pessoas vivendo com HIV (PVHIV), receptores de transplante de órgãos
sólidos ou medula, pacientes em uso de imunobiológicos (especialmente
anti-TNF$\alpha$), contactantes domiciliares de casos bacilíferos, profissionais de
saúde com exposição ocupacional, pacientes em diálise e gestantes em populações de
alta prevalência \cite{who2018, brasil2022}.

O Brasil adota três esquemas terapêuticos prioritários para a ILTB, conforme o
Manual de Recomendações do Ministério da Saúde \cite{brasil2022}: isoniazida diária
por seis ou nove meses (6H ou 9H), rifampicina diária por quatro meses (4R) e
rifapentina + isoniazida semanal por três meses (3HP). A escolha entre esquemas
considera idade, gravidez, função hepática, interações medicamentosas (notadamente
com antirretrovirais em PVHIV) e tolerabilidade, e a literatura clínica
internacional reconhece o 3HP como avanço operacional importante por reduzir o
tempo total de tratamento e melhorar adesão \cite{who2018, who2023}.

A enfermagem desempenha papel central no manejo programático da ILTB no SUS, atuando
no rastreamento ativo de contactantes, na realização da prova tuberculínica (PT) e
do \textit{Interferon-Gamma Release Assay} (IGRA), no acompanhamento mensal de
adesão e na detecção precoce de reações adversas. Esse papel é reconhecido nas
diretrizes da OMS, que identificam o desenvolvimento de recursos humanos e a entrega
de serviços na atenção primária como ações estratégicas para ampliar o tratamento da
ILTB \cite{who2018, who2023}.

Em escala local, \citeauthor{artigo_perfil} \cite{artigo_perfil} caracterizam o
perfil epidemiológico da TB no Espírito Santo entre 2015 e 2020: foram
notificados 8.687 casos no estado no período, com predominância do sexo
masculino (71,06\%) e maior concentração de incidência na faixa etária de 29
a~39 anos. Os municípios da Região Metropolitana de Vitória apresentaram taxas
de incidência substancialmente superiores às do restante do estado, padrão
correlacionado a fatores como tabagismo, etilismo e densidade populacional. Esse
recorte local sustenta a relevância de ferramentas de apoio operacional para os
profissionais da atenção primária do SUS atuantes nesses municípios.

\section{Grandes Modelos de Linguagem (LLMs)}

Os LLMs são modelos de linguagem pré-treinados em corpora massivos, capazes de
realizar diversas tarefas de processamento de linguagem natural com poucos exemplos
\cite{brown2020}. A arquitetura Transformer permite capturar dependências de longo
alcance por meio do mecanismo de atenção.

Para o português, iniciativas como Sabiá-2 \cite{sabia2_2024} e estudos de
adaptação de modelos para o contexto médico brasileiro \cite{adaptingLLMsPortuguese2024,
teachingLLMsBrazil2026} avançaram o estado da arte. Avaliações em exames médicos
brasileiros demonstram desempenho variável entre os modelos disponíveis
\cite{llmsMedicalExam2026}.

Modelos \textit{open-weight} --- com pesos publicamente liberados, ainda que sob
licenças que podem restringir uso comercial --- têm papel particular em contextos
regulatórios estritos. A família Llama, descrita por \citeauthor{dubey2024}
\cite{dubey2024}, exemplifica essa categoria: o release Llama~3 incluiu modelos de
8B a 405B parâmetros, treinados com dados multilíngues e otimizados para
\textit{Grouped Query Attention} (GQA), o que viabiliza inferência eficiente em
\textit{hardware} acessível. O Llama~3.3 70B, refresh subsequente da mesma família
adotado neste trabalho, mantém a capacidade de raciocínio do 405B com fração do
custo computacional, segundo a documentação técnica da Meta.

A relevância prática do paradigma \textit{open-weight} para projetos de saúde em
contextos regulatórios estritos é dupla: \textbf{(i)} permite, em princípio, a
inferência local em servidor institucional, eliminando a transferência de prompts
clínicos a operadores externos --- caminho aderente à LGPD e à recomendação técnica
de evitar envio de dados sensíveis a APIs proprietárias \cite{lgpdSaude2023,
privacyRAGHealthcare2025}; \textbf{(ii)} desacopla a evolução da aplicação do
\textit{roadmap} comercial dos provedores fechados, reduzindo risco de
descontinuidade. No presente projeto, a inferência local não foi adotada por
limitação de \textit{hardware}; a alternativa escolhida foi acessar o Llama~3.3 70B
via API Groq, que não armazena prompts segundo sua política declarada. A discussão
das implicações dessa escolha para implantação institucional consta no
Capítulo~\ref{cap:discussao}.

\section{Retrieval-Augmented Generation (RAG)}

A arquitetura RAG \cite{lewis2021} combina um componente de recuperação de informação
com um LLM generativo. Em vez de depender apenas do conhecimento paramétrico do
modelo, o RAG recupera trechos de uma base de documentos externa e os inclui no
contexto da geração. Isso reduz alucinações e permite atualizar o corpus sem
retreinamento.

\citeauthor{gao2024} \cite{gao2024} categorizam variantes de RAG em básico,
avançado e modular. Estratégias de \textit{chunking} influenciam diretamente a
qualidade da recuperação \cite{advancedChunkingRAG2025}. Técnicas de busca híbrida,
como RRF \cite{cormack2009}, combinam recuperação esparsa (BM25, SPLADE
\cite{formal2021}) com recuperação densa para melhorar o \textit{recall}.

No domínio de saúde, aplicações RAG precisam considerar privacidade e conformidade
regulatória \cite{privacyRAGHealthcare2025, privacyEHRLLMs2025}.

\section{Embeddings e Busca Semântica}

Embeddings são representações vetoriais densas que capturam significado semântico
de textos \cite{boykis2024}. Modelos como Sentence-BERT \cite{reimers2019} permitem
calcular similaridade semântica com eficiência. A busca por vizinhos aproximados
(ANN) com índice HNSW \cite{malkov2016} viabiliza a recuperação em grandes bases.

Para o português clínico, o BioBERTpt \cite{biobertpt2020} foi um dos primeiros
modelos especializados em NLP biomédico, com bons resultados em reconhecimento de
entidades clínicas \cite{clinicalNERPortuguese2026}. O benchmark MMTEB
\cite{mmteb2025} e estudos sobre embeddings para busca clínica
\cite{generalistEmbeddingsClinical2024} fornecem panorama comparativo dos modelos
disponíveis.

Neste trabalho adotou-se o modelo \texttt{paraphrase-multilingual-MiniLM-L12-v2}
(384 dimensões), descendente distilado da família SBERT \cite{reimers2019} treinado
para similaridade semântica multilíngue em mais de 50 idiomas, incluindo o português.
A escolha foi orientada por três critérios: (i) execução local sem dependência de APIs
externas, requisito da política de privacidade do projeto (Seção~\ref{sec:lgpd-ref});
(ii) baixo \textit{footprint} de memória e latência de inferência compatíveis com a
infraestrutura disponível; e (iii) ampla adoção do modelo em sistemas RAG biomédicos
recentes \cite{ocaf2024}.

A literatura, contudo, revela uma tensão entre conveniência operacional e desempenho de
recuperação. Avaliações em corpora hospitalares mostram que modelos de 1024 dimensões,
como \texttt{multilingual-e5-large} e \texttt{BGE-M3}, atingem acurácia
\textit{Top}-10 superior a 90\%, enquanto variantes da família MiniLM permanecem abaixo
de 40\% nos mesmos cenários \cite{SHTI2023, mmteb2025}. Para o português clínico
especificamente, o benchmark SemClinBr \cite{clinicalNERPortuguese2026} indica que
modelos multilíngues genéricos como o mmBERT \cite{marone2025mmbert} superam alternativas monolíngues
especializadas (BioBERTpt, BERTimbau) em reconhecimento de entidades, sugerindo que a
capacidade representacional geral pesa mais que a adaptação de domínio estreita
\cite{generalistEmbeddingsClinical2024}. Esses achados respaldam a estratégia de
empregar um modelo generalista no presente projeto e, ao mesmo tempo, identificam a
migração para um modelo de maior dimensionalidade como vetor natural de evolução
(discutido como trabalho futuro no Capítulo~\ref{cap:conclusao}).

\section{Avaliação de Sistemas RAG}

O \textit{framework} RAGAS \cite{es2024} propõe métricas automáticas para avaliar
sistemas RAG sem anotação humana extensiva:

\begin{itemize}
  \item \textbf{\textit{Faithfulness}}: proporção de afirmações na resposta
        suportadas pelo contexto recuperado;
  \item \textbf{\textit{Answer Relevancy}}: relevância da resposta à pergunta
        original;
  \item \textbf{\textit{Context Precision}}: proporção de trechos recuperados
        úteis para a resposta;
  \item \textbf{\textit{Context Recall}}: cobertura do \textit{ground-truth}
        pelo contexto recuperado.
\end{itemize}

O RAGAS apresenta limitações conhecidas que devem ser ponderadas na interpretação dos
resultados. Em primeiro lugar, todas as métricas dependem de um LLM-juiz para decompor
respostas em afirmações atômicas e verificar suporte no contexto, o que introduz
variabilidade entre execuções e sensibilidade ao modelo avaliador escolhido
\cite{es2024}. Em segundo lugar, a métrica de \textit{Faithfulness} privilegia
respostas extrativas: sínteses parafraseadas, ainda que clinicamente corretas, tendem a
ser penalizadas porque o juiz não consegue mapear cada afirmação literalmente ao
contexto recuperado \cite{bang2023}.

A literatura sobre alucinação em LLMs clínicos diferencia categorias relevantes para a
análise diagnóstica desses scores: alucinações \textit{intrínsecas} (saída contradiz a
fonte), alucinações \textit{extrínsecas} (saída não verificável a partir da fonte) e
alucinações de \textit{citação} (referências inventadas) \cite{ocaf2024}. Cada classe
exige instrumentos de medida distintos, e revisões sistemáticas recentes recomendam
tratar a acurácia de citação como métrica separada da fidelidade ao contexto, sob pena
de subestimar erros graves de proveniência \cite{privacyRAGHealthcare2025}. Por essas
razões, este trabalho complementa a avaliação automática RAGAS com uma sessão
estruturada de avaliação humana com especialista (Capítulo~\ref{cap:avaliacao}),
seguindo a recomendação de combinar métricas automáticas e humanas em RAG biomédico
\cite{ocaf2024}.

\section{Privacidade de Dados e LGPD}
\label{sec:lgpd-ref}

A LGPD (Lei n\textordmasculine\ 13.709/2018) classifica dados de saúde como sensíveis,
sujeitos a requisitos mais rigorosos de consentimento e segurança \cite{lgpdSaude2023,
lgpdEnfermagem2022}. No contexto de LLMs, riscos de vazamento por memorização de
dados de treinamento ou \textit{prompt injection} são documentados
\cite{sokPrivacyLLM2026, privacyEHRLLMs2025}. A literatura específica sobre RAG em
saúde aprofunda esse cenário, identificando que prompts clínicos transmitidos a APIs
de terceiros configuram texto simples no nó de serviço, vulneráveis a inversão de
caixa-preta a partir de ativações intermediárias e à retenção em logs de longo prazo
\cite{privacyRAGHealthcare2025}.

A enfermagem está submetida a obrigações específicas de sigilo profissional. A
Resolução COFEN n\textordmasculine\ 564/2017 (Código de Ética) e a Resolução COFEN
n\textordmasculine\ 429/2012 (registros em prontuário) impõem dever de cautela no
manuseio de dados sensíveis de pacientes, e o descumprimento da LGPD pode ensejar
responsabilização ética e civil \cite{lgpdEnfermagem2022}. No plano regulatório
brasileiro mais amplo, a Autoridade Nacional de Proteção de Dados (Resolução ANPD
n\textordmasculine\ 4/2023) e a Estratégia de Saúde Digital 2020--2028 do Ministério
da Saúde, instituída pelo Decreto n\textordmasculine\ 11.358/2023, balizam o
tratamento de dados em sistemas digitais de saúde \cite{lgpdSaude2023}.

A literatura técnica documenta um conjunto de técnicas para mitigação dos riscos de
vazamento via prompt: \textit{Privacy Span Removal} (substituição seletiva de trechos
sensíveis), \textit{Instance Obfuscation} (\textit{hashing} dinâmico por sessão),
\textit{scrubbing} assistido por modelos de desidentificação e mecanismos de
\textit{Adaptive Differential Privacy} aplicados em nível de \textit{span}
\cite{sokPrivacyLLM2026, privacyRAGHealthcare2025}. Tais técnicas tornam-se
indispensáveis quando o pipeline envia queries clínicas a serviços externos.

No presente projeto, optou-se por uma estratégia conservadora que dispensa essas
mitigações em tempo de execução: (i) o corpus é composto exclusivamente por documentos
públicos do Ministério da Saúde e da OMS, sem dados de paciente; (ii) o pipeline de
embeddings executa-se 100\% local, de modo que queries clínicas não transitam por APIs
de terceiros antes da busca vetorial; (iii) o LLM é acessado via API Groq apenas para
geração da resposta final, sem armazenamento de conversas pelo provedor; e (iv)
nenhuma informação identificável de paciente é inserida nas perguntas de avaliação.
Para a sessão estruturada com o especialista (Capítulo~\ref{cap:avaliacao}), serão
adotados Termo de Consentimento Livre e Esclarecido, política de retenção zero do
histórico de conversa e \textit{logging} mínimo, conforme recomendado pela literatura
para provas de conceito acadêmicas em saúde \cite{privacyRAGHealthcare2025}.

\section{Adoção de IA por Profissionais de Saúde}
\label{sec:adocao-ref}

O Modelo de Aceitação de Tecnologia (TAM) é amplamente usado para explicar a adoção
de sistemas de TI por profissionais de saúde \cite{nursingTAM2024}. Fatores como
utilidade percebida, facilidade de uso e confiança no sistema são determinantes
para a adoção de chatbots clínicos por enfermeiros.

Estudos recentes propõem o TAM-AIN (\textit{Technology Acceptance Model for AI in
Nursing}) como extensão específica para a categoria, incorporando quatro dimensões
ausentes do TAM clássico: (i) \textit{alinhamento ético}, que avalia a compatibilidade
da ferramenta com os valores da profissão e o sigilo do paciente; (ii) \textit{prontidão
organizacional}, que examina o apoio da liderança e a infraestrutura local;
(iii) \textit{preservação da identidade profissional}, relacionada à percepção de
ameaça ou reforço à autonomia clínica e ao componente humano do cuidado; e
(iv) \textit{capacidade de infraestrutura técnica}, que considera as restrições
materiais do ambiente assistencial \cite{nursingTAM2024}. Essa especialização é
particularmente relevante para o SUS, onde a heterogeneidade de infraestrutura entre
unidades de atenção primária pode invalidar premissas de modelos genéricos derivados
de contextos hospitalares de alta complexidade.

A literatura sobre percepção de enfermeiros frente à IA indica padrão consistente de
``otimismo cauteloso'': cerca de 85\% dos profissionais reconhecem benefícios para a
precisão diagnóstica e a redução de erros, mas aproximadamente 70\% mantêm reservas
quanto à preservação do componente humano do cuidado e à segurança dos dados
\cite{nursingTAM2024}. As barreiras mais reportadas incluem infraestrutura tecnológica
obsoleta, sobrecarga de trabalho, ceticismo sobre a capacidade da IA de replicar o
julgamento clínico e, sobretudo entre profissionais com mais tempo de prática, receio
de desqualificação. Esses achados orientam a estrutura da sessão de avaliação com o
especialista (Capítulo~\ref{cap:avaliacao}), que combina rubrica clínica
multidimensional com instrumento TAM-AIN para capturar tanto a acurácia técnica quanto
as dimensões éticas e organizacionais da aceitação.

% ─────────────────────────────────────────────────────────────────────────────
\chapter{Metodologia}
\label{cap:metodologia}
% ─────────────────────────────────────────────────────────────────────────────

Este trabalho adota a \textit{Design Science Research Methodology} (DSRM)
\cite{peffers2007}, metodologia consagrada para pesquisa em sistemas de informação
orientada à construção de artefatos. O DSRM estrutura a pesquisa em seis atividades
iterativas, conforme a Figura~\ref{fig:dsrm}.

\begin{figure}[h]
  \centering
  \begin{minipage}{0.9\textwidth}
    \caption{Processo DSRM aplicado ao desenvolvimento do chatbot ILTB}
    \label{fig:dsrm}
    % \includegraphics[width=\textwidth]{fig/dsrm}
    \aluno{Inserir diagrama DSRM com 6 etapas e iterações}
    \caption*{Fonte: Adaptado de \citeauthor{peffers2007} (\citeyear{peffers2007}).}
  \end{minipage}
\end{figure}

\section{Identificação do Problema e Motivação}

O problema central identificado é a dificuldade de acesso rápido a informações
clínicas padronizadas sobre ILTB por profissionais de enfermagem durante o
atendimento. A solução proposta é um chatbot de apoio à decisão clínica baseado em
RAG.

A identificação do problema partiu de três fontes complementares.

\textbf{Análise documental} dos protocolos vigentes evidenciou que o Manual de
Recomendações do Ministério da Saúde \cite{brasil2022} é extenso, denso em tabelas
e atualizado periodicamente, dificultando a consulta rápida durante o
atendimento. As diretrizes da OMS \cite{who2018, who2023} complementam esse
material e contêm recomendações ocasionalmente desatualizadas em relação às
diretrizes nacionais, exigindo do profissional discernimento sobre qual fonte
prevalece.

\textbf{Revisão da literatura} sobre RAG biomédico \cite{ocaf2024, lewis2021,
gao2024} e sobre adoção de IA por profissionais de enfermagem \cite{nursingTAM2024}
confirmou que ferramentas de apoio à decisão clínica em domínio específico ainda
são pouco exploradas para o português brasileiro, configurando lacuna acadêmica
relevante.

\textbf{Diálogo inicial com profissionais e pesquisadores} da área de
\textit{tuberculose} no contexto de saúde pública --- ainda que tenha tido escopo
reduzido em função das alterações descritas no Capítulo~\ref{cap:introducao} ---
corroborou a percepção de que enfermeiros da atenção primária constituem público
prioritário para uma ferramenta dessa natureza.

A partir dessas três fontes definiu-se a tese central deste trabalho: um chatbot
RAG capaz de responder a perguntas sobre o manejo da ILTB, baseado exclusivamente
em documentos oficiais, em português, com pipeline conforme à LGPD, pode oferecer
apoio operacional relevante a profissionais de enfermagem.

\section{Definição dos Objetivos da Solução}

Os objetivos de design foram definidos com base nos requisitos funcionais e
não funcionais do sistema, derivados da análise do problema. O
Quadro~\ref{quad:decisoes} sintetiza as treze decisões arquiteturais e
metodológicas (D1--D13), com a justificativa bibliográfica de cada uma. Decisões
revisadas no decorrer do projeto --- particularmente D10 e D12, em função do
reescopo da etapa de avaliação --- estão refletidas no quadro em sua versão final.

\begin{quadro}[h]
  \centering
  \caption{Decisões de design do chatbot ILTB (D1--D13)}
  \label{quad:decisoes}
  \begin{tabularx}{\textwidth}{|c|X|X|}
  \hline
  \rowcolor{gray!25}
  \textbf{ID} & \textbf{Decisão} & \textbf{Justificativa} \\ \hline
  D1  & Adoção da arquitetura RAG (recuperação + geração) em vez de LLM puro ou
        \textit{fine-tuning}. & Reduz alucinação em domínio específico ao ancorar
        a geração em documentos oficiais \cite{lewis2021, gao2024, ocaf2024}.
        \\ \hline
  D2  & Embedding multilíngue \texttt{paraphrase-multilingual-MiniLM-L12-v2}
        (384D), executado localmente. & Equilíbrio entre latência,
        \textit{footprint} e requisito LGPD; limitações de dimensionalidade
        documentadas como trabalho futuro \cite{reimers2019, mmteb2025, SHTI2023}.
        \\ \hline
  D3  & \textit{Vector store} ChromaDB com índice HNSW e persistência em disco.
      & Execução embarcada sem servidor dedicado; HNSW provê busca aproximada
        com complexidade logarítmica \cite{malkov2016}. \\ \hline
  D4  & LLM Llama~3.3~70B via API Groq como gerador. & Modelo \textit{open-weight}
        competitivo sem retenção de \textit{prompts} pelo provedor
        \cite{dubey2024}; modelos generalistas grandes superam especializados em
        busca clínica curta \cite{ocaf2024}. \\ \hline
  D5  & \textit{Chunking} semântico por cabeçalhos Markdown, com
        \textit{buffering} de seções pequenas. & Preserva integridade lógica dos
        protocolos; alinhado a recomendações de chunking estrutural
        \cite{gao2024, advancedChunkingRAG2025}. \\ \hline
  D6  & Avaliação automatizada com RAGAS (\textit{faithfulness},
        \textit{answer relevancy}, \textit{context precision/recall}). &
        \textit{Framework reference-free} padrão em RAG biomédico
        \cite{es2024, ocaf2024}. \\ \hline
  D7  & \textit{Pipeline} de embeddings 100\% local; LLM acessado sem
        armazenamento de \textit{prompts}. & Conformidade com LGPD (dados de
        saúde como sensíveis); literatura recomenda evitar APIs externas para
        dados clínicos \cite{lgpdSaude2023, lgpdEnfermagem2022,
        privacyRAGHealthcare2025, sokPrivacyLLM2026}. \\ \hline
  D8  & Adoção da metodologia DSRM como \textit{framework} de pesquisa. & Padrão
        consolidado para pesquisa em sistemas de informação com construção de
        artefato \cite{peffers2007}. \\ \hline
  D9  & Foco em ILTB como problema clínico, com corpus de documentos oficiais
        do MS e OMS. & Doença com alta prevalência global e brasileira; controle
        da ILTB é estratégico para eliminação da TB \cite{who2018, who2023,
        brasil2022}. \\ \hline
  D10 & Avaliação com 1 especialista em ILTB usando TAM-AIN + rubrica clínica
        multidimensional. & Replanejamento decorrente do reescopo do piloto;
        TAM-AIN é instrumento específico para enfermagem \cite{nursingTAM2024,
        ocaf2024}. \\ \hline
  D11 & Busca híbrida (RRF + SPLADE ou ColBERT) registrada como trabalho futuro.
      & \textit{Embedding} denso isolado limita \textit{recall} para termos
        técnicos específicos; literatura sustenta ganhos de busca híbrida
        \cite{cormack2009, formal2021, khattab2020}. \\ \hline
  D12 & Métricas mistas (acurácia vs. MS, severidade Likert, tríade RAG,
        TAM-AIN) em vez de Kappa de Cohen. & Kappa exige $\geq 2$ avaliadores;
        literatura recomenda métricas multidimensionais para validação com $N=1$
        \cite{ocaf2024, es2024}. \\ \hline
  D13 & Documentação explícita da taxa de alucinação como propriedade conhecida
        do LLM, com ablação contra modelo de referência registrada como
        trabalho futuro. & Tipologia e taxas de alucinação reportadas em
        literatura sustentam o teto observado de \textit{faithfulness} no
        Llama~3.3~70B \cite{bang2023, ocaf2024, llmsMedicalExam2026}. \\ \hline
  \end{tabularx}
  \caption*{Fonte: Produção do próprio autor.}
\end{quadro}

\section{Design e Desenvolvimento}

O artefato é um sistema RAG composto pelos seguintes componentes principais:
\textbf{(a)} corpus de documentos clínicos públicos sobre ILTB em português;
\textbf{(b)} pipeline de ingestão (extração, \textit{chunking}, embeddings);
\textbf{(c)} banco de vetores ChromaDB com índice HNSW;
\textbf{(d)} LLM Llama~3.3~70B via API Groq;
\textbf{(e)} API REST (FastAPI) e interface web de demonstração.

\section{Demonstração}

O artefato foi demonstrado via interface web de chat, servida pelo FastAPI e
acessível por meio de túnel ngrok. A demonstração incluiu 38 questões clínicas
sobre ILTB formuladas pelo autor com base nos documentos do corpus.

\section{Avaliação}

A avaliação foi estruturada em duas etapas complementares:

\begin{enumerate}[label=\textbf{(\arabic*)}]
  \item \textbf{Avaliação automatizada (RAGAS):} 38 questões clínicas com
        \textit{ground-truths} extraídos dos documentos indexados, usando
        \textit{faithfulness}, \textit{answer relevancy}, \textit{context
        precision} e \textit{context recall} \cite{es2024};
  \item \textbf{Avaliação qualitativa com especialista:} sessão estruturada com
        profissional especialista em ILTB, avaliando amostras de respostas do
        chatbot em termos de acurácia clínica e aderência aos protocolos vigentes.
\end{enumerate}

A descrição completa do protocolo da sessão com o especialista, incluindo perfil
do avaliador, etapas, instrumentos e rubrica multidimensional, é apresentada na
Seção~\ref{sec:sessao-expert} do Capítulo~\ref{cap:avaliacao}, evitando duplicação.

\section{Comunicação}

Os resultados são comunicados por meio desta monografia e disponibilizados em
repositório público acessível em \url{https://github.com/yago-novaes/chatbot-iltb}.

% ─────────────────────────────────────────────────────────────────────────────
\chapter{Desenvolvimento do Artefato}
\label{cap:desenvolvimento}
% ─────────────────────────────────────────────────────────────────────────────

\section{Arquitetura do Sistema}

O chatbot ILTB é implementado como uma aplicação web com \textit{backend} em
Python (FastAPI) e banco de vetores ChromaDB \cite{chroma2024}. A
Figura~\ref{fig:arquitetura} apresenta a visão geral da arquitetura.

\begin{figure}[h]
  \centering
  \begin{minipage}{0.9\textwidth}
    \caption{Arquitetura do chatbot ILTB baseado em RAG}
    \label{fig:arquitetura}
    % \includegraphics[width=\textwidth]{fig/arquitetura}
    \aluno{Inserir diagrama: usuário → interface web → FastAPI → retriever →
    ChromaDB / LLM via Groq}
    \caption*{Fonte: Produção do próprio autor.}
  \end{minipage}
\end{figure}

\section{Corpus e Pré-processamento}

O corpus indexado é composto por seis documentos públicos sobre o manejo da
tuberculose e da ILTB, com total de 551~páginas, conforme detalhado na
Tabela~\ref{tab:corpus}. A seleção integra três fontes ministeriais brasileiras
(MS), uma da OPAS/OMS em português e uma de sociedade médica nacional (GEDIIB),
garantindo: (i) cobertura completa do protocolo brasileiro vigente; (ii)
alinhamento com recomendações internacionais; e (iii) domínio público ---
requisito para a estratégia de privacidade adotada (Seção~\ref{sec:lgpd-ref}). A
sustentação teórica do referencial clínico (Seção~\ref{cap:referencial}) apoia-se
adicionalmente nos documentos canônicos da OMS \cite{who2018, who2023} e no
\textit{Manual de Recomendações para o Controle da Tuberculose no Brasil}
\cite{brasil2022}.

\begin{table}[h]
  \centering
  \caption{Documentos indexados no corpus do chatbot ILTB}
  \label{tab:corpus}
  \begin{tabularx}{\textwidth}{@{}lXrr@{}}
  \toprule
  \textbf{Identificador} & \textbf{Título} & \textbf{Ano} & \textbf{Pág.} \\
  \midrule
  \textsc{ms-manual}       & Manual de Recomendações para o Controle da
                             Tuberculose no Brasil (MS, 2ª ed.\ atualizada) &
                             2019 & 366 \\
  \textsc{ms-guia}         & Recomendações para o Controle da Tuberculose ---
                             Guia Rápido para Profissionais de Saúde (MS, 2ª
                             ed.) & 2021 & 50 \\
  \textsc{ms-vig-iltb}     & Protocolo de Vigilância da Infecção Latente pelo
                             \textit{M.\,tuberculosis} no Brasil (MS, 2ª ed.) &
                             2022 & 38 \\
  \textsc{ms-3hp}          & Tratamento da Infecção Latente com Rifapentina +
                             Isoniazida (3HP) --- MS/SVS & 2021 & 2 \\
  \textsc{opas-mod4}       & Manual Operacional de Tuberculose da OMS ---
                             Módulo 4: Tratamento (OPAS/OMS, versão em
                             português) & 2020 & 84 \\
  \textsc{gediib-tb}       & Tratamento da Tuberculose Infecção: Posicionamentos
                             Científicos do GEDIIB & 2024 & 11 \\
  \midrule
  \multicolumn{3}{l}{\textbf{Total}} & \textbf{551} \\
  \bottomrule
  \end{tabularx}
  \caption*{Fonte: Produção do próprio autor.}
\end{table}

O \textit{pipeline} de ingestão extrai texto e estrutura dos PDFs em duas etapas:
extração inicial com \texttt{pypdf}, seguida por sanitização e reorganização em
Markdown estruturado por cabeçalhos (\texttt{\#}~a~\texttt{\#\#\#\#}). Durante a
extração foram identificados desafios típicos de PDFs governamentais
brasileiros, com destaque para a desconfiguração de tabelas (esquemas
terapêuticos, doses por faixa etária) e uma falha de \textit{parsing} que produzia
\textit{gap} de 34 páginas no Manual do MS. Esse defeito foi diagnosticado e
corrigido por meio de \textit{patches} pontuais no \textit{stream} extraído,
documentados no diário técnico do projeto.

A sanitização Markdown padroniza unidades, doses e nomes de fármacos, removendo
artefatos de \textit{OCR} e normalizando hierarquia de seções. O produto final
é um conjunto de arquivos \texttt{.md} estruturados, prontos para a etapa de
\textit{chunking}.

O conjunto totaliza 551~páginas (Tabela~\ref{tab:corpus}), e após a sanitização
e o \textit{chunking} (Seção~\ref{cap:desenvolvimento}) é convertido em 898
\textit{chunks} indexados, distribuídos majoritariamente entre o
\textsc{ms-manual} (que concentra a maior densidade de protocolo clínico) e o
\textsc{opas-mod4}.

\section{Segmentação (\textit{Chunking}) e Indexação}

A estratégia de \textit{chunking} adotada é hierárquica e dirigida por estrutura:
o documento é segmentado nas fronteiras de cabeçalho Markdown
(\texttt{\#} a \texttt{\#\#\#\#}), reproduzindo a hierarquia conceitual dos
protocolos. Essa abordagem alinha-se às recomendações de \citeauthor{gao2024}
\cite{gao2024} de privilegiar fronteiras semânticas naturais sobre janelas fixas,
e às evidências quantitativas reportadas para \textit{chunking} estrutural em
domínios técnicos \cite{advancedChunkingRAG2025}.

Para evitar fragmentação excessiva de seções curtas, o \textit{chunker} aplica
\textit{buffering}: seções abaixo de um tamanho mínimo
(\texttt{MIN\_CHUNK\_SIZE}~$=400$ caracteres) são agregadas à seção subsequente do
mesmo nível hierárquico. Seções acima do limite superior são subdivididas por
parágrafos, preservando coerência sintática. A regex de detecção de cabeçalho é
\texttt{\textasciicircum\#\{1,4\}\textbackslash s} para capturar os quatro níveis
relevantes no corpus.

O processo de \textit{chunking} foi iterado ao longo da Fase~2 do desenvolvimento,
com avaliações RAGAS sucessivas para calibrar o parâmetro \texttt{MIN\_CHUNK\_SIZE}.
A configuração final produziu \textbf{898 \textit{chunks}} a partir dos três
documentos indexados, distribuídos majoritariamente entre o Manual do MS
\cite{brasil2022} e as diretrizes da OMS \cite{who2018}. Esses \textit{chunks}
são persistidos no ChromaDB com seus respectivos vetores e metadados de origem
(arquivo, hierarquia de cabeçalhos, posição relativa).

\section{Modelo de Embeddings}

O modelo de \textit{embeddings} adotado é o
\texttt{paraphrase-multilingual-MiniLM-L12-v2}, da família SBERT
\cite{reimers2019}, com 384 dimensões e suporte a mais de 50 idiomas. Foi carregado
via biblioteca \texttt{sentence-transformers} e executado localmente, sem requisição
a APIs externas, como descrito conceitualmente na
Seção~\ref{cap:referencial} e no item D2 do Quadro~\ref{quad:decisoes}.

A escolha foi orientada por três critérios operacionais já discutidos: (i)
conformidade com a política LGPD do projeto (Seção~\ref{sec:lgpd-ref}); (ii)
\textit{footprint} compatível com a infraestrutura disponível; (iii) ampla
adoção em RAG biomédico \cite{ocaf2024}. O modelo gera vetores de 384 dimensões
que são indexados no ChromaDB com métrica de similaridade cossenoidal.

\textbf{Limitação observada empiricamente.} Um experimento conduzido durante o
desenvolvimento adicionou contexto hierárquico explícito a cada \textit{chunk}
(prefixo com a cadeia de cabeçalhos), prática conhecida como
\textit{contextual chunking}. Essa modificação \textit{piorou} os resultados
RAGAS em todas as métricas, comportamento consistente com a literatura
\cite{SHTI2023}: modelos de baixa dimensionalidade (384D) misturam semântica
estrutural com semântica clínica no vetor resultante, degradando a precisão.
Modelos de maior dimensionalidade ($\geq 1024$D) absorvem melhor essa
informação adicional, o que reforça a indicação de migração para um modelo
1024D (e.g., \texttt{multilingual-e5-large}, \texttt{BGE-M3}) como trabalho
futuro prioritário (Capítulo~\ref{cap:conclusao}).

\section{Módulo de Recuperação}

O módulo de recuperação utiliza busca por similaridade cossenoidal sobre o índice
HNSW do ChromaDB. Dado um \textit{query} do usuário, os $k$ trechos mais similares
são recuperados e concatenados como contexto para o LLM.

Após calibração experimental por meio de execuções RAGAS, o parâmetro $k$ foi
fixado em \textbf{5}. Valores menores ($k=3, 4$) excluíam, com frequência,
trechos relevantes em consultas envolvendo populações especiais ou interações
medicamentosas; valores maiores ($k \geq 7$) introduziam ruído sem ganho de
\textit{recall}.

Adicionalmente, foi adotado um \textbf{limiar de similaridade cossenoidal de
$0{,}40$}. Consultas cujo \textit{top-1} apresenta \textit{score} abaixo desse
limiar acionam um \textit{fallback} que orienta o modelo a responder
explicitamente que a informação não foi encontrada no corpus, em lugar de
sintetizar uma resposta apoiada em fragmentos pouco relevantes. Essa salvaguarda
mitiga o risco de alucinação extrínseca em consultas fora do escopo dos
protocolos indexados.

O limiar foi definido empiricamente a partir da distribuição de \textit{scores}
nas 38 questões de avaliação: valores acima de $0{,}40$ correlacionaram-se com
\textit{ground truths} extraíveis dos \textit{chunks} recuperados, enquanto
valores abaixo coincidiam, em sua maioria, com tópicos fora do corpus.

\section{Módulo de Geração}

O modelo Llama~3.3~70B \cite{dubey2024} é acessado via API Groq. O \textit{prompt
system} instrui o modelo a responder exclusivamente com base no contexto recuperado,
em português do Brasil, e a indicar quando a informação não está disponível no
corpus.

O \textit{system prompt} de produção (versão~v1) instrui o modelo a: (i)
responder \textbf{exclusivamente} com base no contexto recuperado; (ii) usar
português brasileiro com terminologia clínica adequada à enfermagem; (iii)
declarar explicitamente quando a informação não estiver no contexto
(comportamento alinhado ao \textit{fallback} descrito na seção anterior);
(iv) preservar números, doses e unidades exatamente como no protocolo
original; e (v) anexar referência à fonte (documento de origem) quando
disponível nos metadados. O texto integral do \textit{prompt} está reproduzido
no Apêndice~\ref{ape:prompt}.

\section{Interface de Demonstração}

A interface de demonstração é uma aplicação web responsiva (HTML/CSS/JavaScript)
servida pelo FastAPI e acessível via túnel ngrok. A interface inclui aviso explícito
de que a ferramenta é de apoio e não substitui o julgamento clínico profissional.

\aluno{Inserir screenshot da interface de chat.}

% ─────────────────────────────────────────────────────────────────────────────
\chapter{Avaliação e Resultados}
\label{cap:avaliacao}
% ─────────────────────────────────────────────────────────────────────────────

\section{Dataset de Avaliação}

O \textit{dataset} de avaliação RAGAS foi construído pelo autor a partir do corpus
indexado, totalizando \textbf{42 questões clínicas} sobre o manejo da ILTB, das
quais \textbf{38} são \textit{in-scope} (utilizadas no cálculo das métricas) e
\textbf{4} são \textit{out-of-scope}, destinadas exclusivamente ao teste do
\textit{fallback} de similaridade descrito na Seção~\ref{cap:desenvolvimento}.

As 38 questões \textit{in-scope} estão organizadas em sete categorias temáticas
(Tabela~\ref{tab:dataset-cat}), que cobrem os principais eixos de consulta
esperados em uso real.

\begin{table}[h]
  \centering
  \caption{Distribuição do \textit{dataset} de avaliação por categoria}
  \label{tab:dataset-cat}
  \begin{tabular}{@{}clc@{}}
  \toprule
  \textbf{Sigla} & \textbf{Categoria} & \textbf{N} \\
  \midrule
  ET   & Esquemas terapêuticos (6H, 9H, 4R, 3HP)            & 7 \\
  IND  & Indicações de tratamento                            & 5 \\
  PE   & Populações especiais (PVHIV, gestantes, crianças)   & 7 \\
  DI   & Diagnóstico (PT, IGRA, exclusão de TB ativa)        & 5 \\
  MO   & Monitoramento clínico e laboratorial                & 5 \\
  IT   & Interações medicamentosas                           & 5 \\
  EA   & Efeitos adversos                                    & 4 \\
  \midrule
  \multicolumn{2}{l}{\textit{Total \textit{in-scope}}}        & \textbf{38} \\
  FE   & Fora do escopo (apenas para teste do \textit{fallback}) & 4 \\
  \midrule
  \multicolumn{2}{l}{\textit{Total geral}}                     & \textbf{42} \\
  \bottomrule
  \end{tabular}
  \caption*{Fonte: Produção do próprio autor. Dataset completo no
            Apêndice~\ref{ape:dataset}.}
\end{table}

Os \textit{ground truths} foram extraídos \textbf{literalmente} dos arquivos
Markdown indexados, preservando trechos verbatim do MS ou da OMS. Essa decisão
metodológica fortalece a interpretação das métricas RAGAS: como o \textit{ground
truth} é, por construção, recuperável do corpus, eventuais falhas no
\textit{Context Recall} indicam problemas reais de recuperação --- não ausência
de informação na base.

Quatro questões originalmente incluídas tratavam de tópicos \textbf{fora do
escopo} (\textit{out-of-scope}), destinadas a testar o \textit{fallback} de
similaridade descrito na Seção~\ref{cap:desenvolvimento}; elas são excluídas dos
cálculos RAGAS, totalizando $36$--$38$ amostras válidas por execução conforme
disponibilidade do LLM-juiz.

O LLM-juiz adotado para a avaliação foi o \texttt{gpt-4o-mini} (OpenAI),
escolhido por três motivos: (i) suporte a $n>1$ \textit{completions}, necessário
para o cálculo correto de \textit{Faithfulness}; (ii) ausência de limite TPM
restritivo que inviabilize execuções completas; (iii) custo aproximado de
US\$\,$0{,}05$ por execução de 38 questões. Importante notar que o
\textit{gpt-4o-mini} foi usado \textbf{exclusivamente como avaliador},
permanecendo o Llama~3.3~70B via Groq como o gerador no \textit{pipeline} de
produção.

\section{Avaliação Automatizada (RAGAS)}

A Tabela~\ref{tab:ragas} apresenta a evolução dos resultados ao longo das principais
configurações testadas.

\begin{table}[ht]
  \caption{Evolução dos resultados RAGAS do chatbot ILTB}
  \label{tab:ragas}
  \begin{tabularx}{\textwidth}{@{}Xcccc@{}}
  \toprule
  \textbf{Configuração} &
  \textbf{Faith.} &
  \textbf{Ans. Rel.} &
  \textbf{Ctx. Prec.} &
  \textbf{Ctx. Rec.} \\
  \midrule
  Baseline (pré-sanitização, top\_k=4) & 0{,}375 & 0{,}310 & 0{,}548 & 0{,}382 \\
  Pós-sanitização, prompt v1, top\_k=4 & 0{,}586 & 0{,}534 & 0{,}656 & 0{,}608 \\
  A/B gpt-4o-mini como pipeline LLM     & 0{,}525 & 0{,}524 & 0{,}657 & 0{,}583 \\
  \textbf{Chunker hierárquico, top\_k=5 (final)} &
    \textbf{0{,}515} & \textbf{0{,}381} & \textbf{0{,}735} & \textbf{0{,}520} \\
  \bottomrule
  \end{tabularx}
  \caption*{Fonte: Produção do próprio autor.}
\end{table}

A Tabela~\ref{tab:ragas} sintetiza quatro pontos da trajetória de otimização, dos
quais decorrem as seguintes leituras:

\textbf{Linha 1 (\textit{baseline}).} Configuração inicial sem sanitização de
\textit{chunks} e com $k=4$. Os \textit{scores} baixos em todas as métricas
indicam principalmente falhas de recuperação: \textit{Context Precision} de
$0{,}548$ confirma que aproximadamente metade dos \textit{chunks} recuperados
não eram relevantes para a pergunta, propagando essa falha para
\textit{Faithfulness} e \textit{Answer Relevancy}.

\textbf{Linha 2 (pós-sanitização, \textit{prompt} v1).} A introdução de
sanitização Markdown (Seção~\ref{cap:desenvolvimento}) e a primeira versão
estável do \textit{system prompt} elevaram \textit{Faithfulness} a $0{,}586$,
patamar próximo ao observado em revisões de RAG biomédico \cite{ocaf2024}, e
duplicaram o \textit{Answer Relevancy}. Esse foi o salto qualitativo principal
do projeto.

\textbf{Linha 3 (\textit{gpt-4o-mini} no \textit{pipeline}).} Teste A/B
substituindo o Llama~3.3~70B pelo \texttt{gpt-4o-mini} também como gerador
(não apenas como juiz). Os resultados \textit{não} mostraram superioridade
relevante, e a opção foi descartada, mantendo o Llama por razões de custo e
licença \textit{open-weight}.

\textbf{Linha 4 (encerramento da Fase 2).} Adoção do \textit{chunker} hierárquico
descrito na Seção~\ref{cap:desenvolvimento} e $k=5$. O \textit{Context Precision}
alcançou $0{,}735$ --- próximo à meta de $0{,}75$ ---, indicando que a recuperação
estabilizou em torno de um patamar útil. O \textit{Faithfulness}, em contrapartida,
permaneceu em $0{,}515$, abrindo a questão diagnóstica investigada pela ablação
$2 \times 2$ apresentada na Seção~\ref{sec:ablacao-2x2}.

\subsection{Ablação 2×2: Gerador e Embedding}
\label{sec:ablacao-2x2}

O patamar de \textit{Faithfulness} observado ao fim da Fase 2 admitia duas
explicações concorrentes: (i) limitação do \textbf{gerador} --- o Llama~3.3~70B
seria intrinsecamente insuficiente para sintetizar respostas fiéis ao contexto
no domínio ILTB; ou (ii) limitação do \textit{pipeline} \textbf{RAG} ---
\textit{retriever}/\textit{embedding} insuficientes para entregar contexto rico,
forçando o gerador a respostas pobres. A literatura de RAG clínico recomenda,
exatamente nessa situação, uma ablação cruzada isolando cada variável
\cite{ocaf2024, wang2024bestpractices}.

Realizou-se, portanto, uma ablação $2 \times 2$ sobre o mesmo conjunto de 38
questões \textit{in-scope} (Seção anterior), variando dois fatores enquanto
mantendo o restante do \textit{pipeline} inalterado: o \textbf{gerador}
(Llama~3.3~70B vs GPT-4o) e o \textbf{\textit{embedding}}
(\texttt{paraphrase-multilingual-MiniLM-L12-v2} 384D vs \texttt{BAAI/bge-m3}
1024D). Quatro corridas independentes foram executadas, cada uma com a coleta
das respostas e o cálculo das métricas RAGAS pelo mesmo juiz LLM
(\texttt{gpt-4o-mini}). A Tabela~\ref{tab:ablacao-2x2} apresenta os resultados.

\begin{table}[ht]
  \caption{Ablação $2 \times 2$: Gerador × \textit{Embedding} (38 questões; juiz
           RAGAS \texttt{gpt-4o-mini})}
  \label{tab:ablacao-2x2}
  \begin{tabularx}{\textwidth}{@{}lXcccc@{}}
  \toprule
  \textbf{Gerador} & \textbf{\textit{Embedding}} &
  \textbf{Faith.} & \textbf{Ans. Rel.} & \textbf{Ctx. Prec.} & \textbf{Ctx. Rec.} \\
  \midrule
  Llama~3.3~70B & MiniLM 384D (gate Fase 2)              & 0{,}515 & 0{,}381 & 0{,}735 & 0{,}520 \\
  GPT-4o        & MiniLM 384D                            & 0{,}383 & 0{,}315 & 0{,}761 & 0{,}533 \\
  GPT-4o        & BGE-M3 1024D                           & 0{,}600 & 0{,}618 & 0{,}907 & 0{,}796 \\
  \textbf{Llama~3.3~70B} & \textbf{BGE-M3 1024D (gate atual)}
                                                        & \textbf{0{,}675} & \textbf{0{,}686} & \textbf{0{,}949} & \textbf{0{,}740} \\
  \bottomrule
  \end{tabularx}
  \caption*{Fonte: Produção do próprio autor.}
\end{table}

\textbf{Efeito do \textit{embedding} (linhas pareadas, gerador fixo).}
Mantendo o gerador GPT-4o constante, a substituição MiniLM 384D → BGE-M3 1024D
produziu ganhos de $+56{,}7\%$ em \textit{Faithfulness} ($0{,}383 \to 0{,}600$),
$+96{,}2\%$ em \textit{Answer Relevancy} ($0{,}315 \to 0{,}618$), $+19{,}2\%$ em
\textit{Context Precision} ($0{,}761 \to 0{,}907$, ultrapassando confortavelmente
a meta de $0{,}75$) e $+49{,}3\%$ em \textit{Context Recall}
($0{,}533 \to 0{,}796$). Mantendo o gerador Llama constante, a mesma substituição
produziu ganhos de $+31\%$, $+80\%$, $+29\%$ e $+42\%$, respectivamente. Em ambos
os geradores, portanto, a migração de \textit{embedding} produz salto material em
todas as métricas --- direção e magnitude consistentes com os \textit{benchmarks}
multilíngues \cite{SHTI2023, mmteb2025, marone2025mmbert} e com a revisão
sistemática de \citeauthor{ocaf2024} \cite{ocaf2024}, que identifica a escolha
de \textit{embedding} como o principal lever de qualidade em \textit{pipelines}
RAG clínicos.

\textbf{Efeito do gerador (colunas pareadas, \textit{embedding} fixo).}
Surpreendentemente, substituir o Llama pelo GPT-4o \emph{piora} a
\textit{Faithfulness} em ambos os \textit{embeddings}: $0{,}515 \to 0{,}383$ no
MiniLM ($-25{,}6\%$) e $0{,}675 \to 0{,}600$ no BGE-M3 ($-11{,}1\%$). O
\textit{Context Precision} permanece estatisticamente equivalente entre os dois
geradores (esperado --- essa métrica depende do \textit{retriever}, não do
gerador). \textit{Answer Relevancy} também cai com GPT-4o em ambos os
\textit{embeddings}. A inspeção direta das respostas geradas revelou a causa:
o GPT-4o segue mais estritamente as instruções do \textit{system prompt}
(Apêndice~\ref{ape:prompt}) e produz, com maior frequência, a resposta de
\textit{fallback} literal --- ``Não encontrei essa informação nos protocolos
indexados.'' --- mesmo em casos em que o conteúdo está presente nos
\textit{chunks} recuperados. O Llama, em contraste, sintetiza a partir dos
\textit{chunks} mesmo quando o protocolo de origem é parcial, comportamento que
o RAGAS pontua como mais fiel porque há \textit{claims} verificáveis na resposta.

Esse efeito --- modelo mais alinhado a instruções pontuando \emph{pior} em
métricas automatizadas --- é convergente com observações independentes em RAG
clínico: \citeauthor{zakka2024almanac} \cite{zakka2024almanac}, no sistema
Almanac (NEJM AI), relata que modelos de fronteira frequentemente recusam
responder em contextos clínicos e exigem engenharia específica para liberar a
síntese quando segura. A literatura de mitigação de alucinação em LLMs
\cite{huang2023halucination, ocaf2024} descreve esse padrão como ``abstenção
excessiva'' --- evitar alucinação ao custo de cobertura. O paper original do
RAGAS \cite{es2024} reconhece, igualmente, que sua métrica de
\textit{Faithfulness} pode penalizar respostas seguras-por-recusa em cenários
de síntese.

\subsection{Configuração Pós-Migração (Intermediária)}
\label{sec:gate-pos-migracao}

A célula \textbf{Llama~3.3~70B + BGE-M3 1024D} (Tabela~\ref{tab:ablacao-2x2},
linha 4) constituiu o gate intermediário do artefato após a migração de
\textit{embedding}. A comparação direta com o gate da Fase 2 mostra ganhos
uniformes em todas as métricas: \textit{Faithfulness} $0{,}515 \to 0{,}675$
($+31\%$); \textit{Answer Relevancy} $0{,}381 \to 0{,}686$ ($+80\%$);
\textit{Context Precision} $0{,}735 \to 0{,}949$ ($+29\%$, agora PASS robusto
sobre a meta $\geq 0{,}75$); \textit{Context Recall} $0{,}520 \to 0{,}740$
($+42\%$).

A \textit{Faithfulness} de $0{,}675$ permanece abaixo da meta original de
$0{,}80$, mas o \textbf{diagnóstico do gap} se alterou substancialmente. Com
\textit{Context Precision} em $0{,}949$, o problema deixou de ser de
\textbf{recuperação} e tornou-se de \textbf{geração}: o \textit{retriever}
entrega \textit{chunks} relevantes, mas o gerador ainda produz, em parte das
questões, afirmações que o juiz RAGAS classifica como não estritamente
derivadas do contexto. Esse diagnóstico motivou três intervenções
pós-recuperação subsequentes, descritas nas Seções~\ref{sec:hybrid}
(busca híbrida com BM25), \ref{sec:rerank} (\textit{reranking}
\textit{cross-encoder}) e \ref{sec:chunk-test} (tuning de \textit{chunk\_size}),
fechando a iteração experimental do trabalho.

\subsection{Iteração 1: Busca Híbrida com BM25 e RRF}
\label{sec:hybrid}

A primeira intervenção pós-recuperação adotou busca híbrida combinando o
\textit{retriever} denso BGE-M3 já em produção com um índice esparso BM25
(\citeauthor{cormack2009}, \citeyear{cormack2009}; \citeauthor{formal2021},
\citeyear{formal2021}), fundindo os dois \textit{rankings} via
\textit{Reciprocal Rank Fusion} (RRF). A hipótese é que termos clínicos exatos
(siglas como ``PVHIV'' e ``3HP'', dosagens, nomes de fármacos) e construções
com negação respondem melhor a casamento lexical, complementando o
\textit{retriever} densamente codificado.

O índice BM25 é construído como \textit{sidecar} do ChromaDB --- compartilha
os mesmos 898 \textit{chunks} indexados, sem duplicação. Cada consulta gera
20 candidatos de cada \textit{ranker}, e o RRF (constante $k = 60$, conforme
\citeauthor{cormack2009}) computa o \textit{score} final
$\sum_{r} 1 / (k + \text{rank}_r(d))$, do qual se selecionam os $k = 5$
\textit{chunks} de maior \textit{rank} fundido.

A Tabela~\ref{tab:hybrid-rerank} (linha 2) apresenta o resultado da avaliação
RAGAS sobre as mesmas 38 questões. Os ganhos materiais foram em
\textit{Context Recall} ($+11{,}6\%$) e \textit{Answer Relevancy} ($+7{,}4\%$):
o BM25 captura \textit{chunks} com termos exatos que o BGE-M3 não recuperava,
e o gerador, com cobertura maior, produz respostas mais on-topic. O ganho em
\textit{Faithfulness} foi modesto ($+2{,}1\%$, $0{,}675 \to 0{,}689$), abaixo
do envelope esperado pela literatura
($+3$ a $+7$\,pp, \citeauthor{wang2024bestpractices}, \citeyear{wang2024bestpractices})
--- explicado pelo BGE-M3 já ser muito forte em multilíngue
(\textit{state-of-the-art} no MMTEB, conforme \cite{mmteb2025}), o que reduz
a informação nova trazida pelo BM25. Um efeito secundário relevante foi a
melhoria do comportamento de \textit{fallback} em consultas \textit{out-of-scope}:
\textit{scores} RRF nessas consultas caíram para 0{,}03 (vs. $\sim 0{,}87$
no \textit{retrieval} denso puro), facilitando rejeição segura por
limiar simples.

\subsection{Iteração 2: \textit{Reranking} \textit{Cross-encoder} (mGTE)}
\label{sec:rerank}

Sobre o \textit{pipeline} híbrido da Seção anterior, adicionou-se uma etapa
de \textit{reranking} com o \textit{cross-encoder} multilíngue
\texttt{Alibaba-NLP/gte-multilingual-reranker-base} (\citeauthor{zhang2024mgte},
\citeyear{zhang2024mgte}). \textit{Cross-encoders} processam o par
(\textit{query}, \textit{chunk}) conjuntamente, capturando relevância
fine-grained que \textit{bi-encoders} densos como o BGE-M3 e modelos esparsos
como o BM25 não capturam isoladamente.

O fluxo final é: busca densa top-20 $\cup$ BM25 top-20 $\to$ fusão RRF $\to$
\textit{cross-encoder} reordena $\to$ top-5 enviado ao gerador. O custo
computacional é aceitável (\textasciitilde{}200\,ms adicionais por consulta
em CPU), sem necessidade de GPU dedicada.

A Tabela~\ref{tab:hybrid-rerank} (linha 3) sumariza o resultado. O
\textit{reranking} produziu o salto mais expressivo da iteração:
\textit{Faithfulness} $0{,}689 \to 0{,}761$ ($+10{,}4\%$ sobre o híbrido sem
\textit{rerank}; $+12{,}7\%$ sobre o gate intermediário). Esse ganho recai
exatamente no envelope previsto por
\citeauthor{wang2024bestpractices} ($+5$ a $+10$\,pp). \textit{Context
Precision} também melhorou ($0{,}912 \to 0{,}964$, $+5{,}7\%$) ---
o \textit{cross-encoder} filtra os \textit{chunks} lexicalmente similares mas
semanticamente periféricos que o BM25 ocasionalmente promove. As pequenas
regressões em \textit{Answer Relevancy} ($-2{,}3\%$) e \textit{Context Recall}
($-2{,}8\%$) refletem o comportamento mais conservador do
\textit{reranker}: descarta candidatos marginais que ampliariam a cobertura
mas introduziriam ruído.

\begin{table}[ht]
  \caption{Iterações pós-recuperação aplicadas sobre Llama~3.3~70B + BGE-M3 1024D}
  \label{tab:hybrid-rerank}
  \begin{tabularx}{\textwidth}{@{}Xcccc@{}}
  \toprule
  \textbf{Configuração} & \textbf{Faith.} & \textbf{Ans. Rel.} &
  \textbf{Ctx. Prec.} & \textbf{Ctx. Rec.} \\
  \midrule
  Pós-migração (denso BGE-M3 apenas)         & 0{,}675 & 0{,}686 & 0{,}949 & 0{,}740 \\
  + Busca híbrida (BGE-M3 + BM25, RRF)       & 0{,}689 & 0{,}737 & 0{,}912 & 0{,}826 \\
  \textbf{+ \textit{Reranker} mGTE (gate final)}
                                             & \textbf{0{,}761} & \textbf{0{,}720} & \textbf{0{,}964} & \textbf{0{,}803} \\
  \bottomrule
  \end{tabularx}
  \caption*{Fonte: Produção do próprio autor. 38 questões \textit{in-scope};
            juiz RAGAS \texttt{gpt-4o-mini}.}
\end{table}

\subsection{Iteração 3: Teste de \textit{chunk\_size}}
\label{sec:chunk-test}

A literatura \cite{wang2024bestpractices} sugere que \textit{chunks} menores
(300--500 \textit{tokens}) otimizam o trade-off entre fidelidade e cobertura
em domínios técnicos. Avaliou-se, portanto, a re-ingestão do corpus com
\texttt{CHUNK\_SIZE = 500} caracteres (vs. 800 da configuração padrão), sobre
o \textit{pipeline} \textit{hybrid+rerank} já estabilizado. Resultados sobre
as mesmas 38 questões: \textit{Faithfulness} $0{,}761 \to 0{,}756$ ($-0{,}005$,
ruído de juiz); \textit{Answer Relevancy} $0{,}720 \to 0{,}683$ ($-5{,}1\%$,
regressão material); \textit{Context Precision} $0{,}964 \to 0{,}951$
($-1{,}3\%$); \textit{Context Recall} $0{,}803 \to 0{,}811$ ($+1{,}0\%$).

Diagnóstico do efeito nulo: o \textit{chunker} hierárquico do
\textit{pipeline} (descrito na Seção~\ref{cap:desenvolvimento}) é guiado
pelas fronteiras de cabeçalho do \textit{Markdown} sanitizado e por um piso
\texttt{MIN\_CHUNK\_SIZE = 400} \textit{chars}. Reduzir o teto de 800 para
500 \textit{chars} alterou apenas marginalmente a distribuição de tamanhos
(mediana de \textasciitilde{}800 para \textasciitilde{}674 \textit{chars}),
preservando a maior parte da granularidade original. Para um teste estrito
do efeito descrito por \citeauthor{wang2024bestpractices} seria necessário
substituir o \textit{chunker} por uma estratégia baseada em \textit{tokens}
do modelo (\textit{e.g.}, \textit{Small-to-Big} \cite{advancedChunkingRAG2025}),
mudança fora do escopo desta iteração. O \texttt{CHUNK\_SIZE = 800} foi,
portanto, mantido.

\subsection{Configuração Final do Artefato}
\label{sec:gate-final}

A combinação Llama~3.3~70B + BGE-M3 1024D + busca híbrida (BM25 + denso, RRF
$k = 60$) + \textit{reranker} mGTE foi adotada como \textbf{gate operacional
final do artefato}. A Tabela~\ref{tab:gate-final} consolida os
\textit{scores} comparados ao gate da Fase 2.

\begin{table}[ht]
  \caption{Gate operacional final vs. gate da Fase 2 (\textit{baseline} pré-migração)}
  \label{tab:gate-final}
  \begin{tabularx}{\textwidth}{@{}Xcccc@{}}
  \toprule
  \textbf{Configuração} &
  \textbf{Faith.} & \textbf{Ans. Rel.} &
  \textbf{Ctx. Prec.} & \textbf{Ctx. Rec.} \\
  \midrule
  Gate Fase 2 (\textit{baseline}, MiniLM 384D + denso) & 0{,}515 & 0{,}381 & 0{,}735 & 0{,}520 \\
  \textbf{Gate Final} (BGE-M3 + híbrido + \textit{rerank}) &
    \textbf{0{,}761} & \textbf{0{,}720} & \textbf{0{,}964} & \textbf{0{,}803} \\
  \midrule
  \textit{$\Delta$ absoluto} &
    $+0{,}246$ & $+0{,}339$ & $+0{,}229$ & $+0{,}283$ \\
  \textit{$\Delta$ relativo} &
    $+47{,}8\%$ & $+89{,}0\%$ & $+31{,}2\%$ & $+54{,}4\%$ \\
  \textit{Gate (faith $\geq 0{,}80$; ctx\_prec $\geq 0{,}75$)} &
    FAIL / FAIL & — & \textbf{FAIL / PASS} & — \\
  \bottomrule
  \end{tabularx}
  \caption*{Fonte: Produção do próprio autor.}
\end{table}

\textit{Context Precision} ultrapassa confortavelmente a meta de $0{,}75$
($+0{,}214$ sobre o alvo), indicando que o \textit{retrieval} convergiu para
um patamar maduro. \textit{Faithfulness} de $0{,}761$ está a $0{,}039$ da
meta original de $0{,}80$ --- gap residual interpretado na
Seção~\ref{sec:interp-faith} e endereçado, como linha de pesquisa imediata,
no Capítulo~\ref{cap:conclusao}.

\subsection{Interpretação da Métrica de \textit{Faithfulness}}
\label{sec:interp-faith}

O patamar de \textit{Faithfulness} no gate final do artefato, $0{,}761$
(Tabela~\ref{tab:gate-final}), está acima do que a literatura de RAG clínico
reporta como teto comum sem intervenções de \textit{self-reflection}: a
revisão sistemática de \citeauthor{ocaf2024} \cite{ocaf2024}, combinando 30
estudos do domínio, indica que valores acima de $0{,}75$ raramente são
atingidos sem mecanismos adicionais sobre o esquema
\textit{retrieve-then-generate} básico. As três iterações pós-recuperação
descritas nas Seções \ref{sec:hybrid}--\ref{sec:gate-final} (busca híbrida,
\textit{reranking} \textit{cross-encoder} e teste de \textit{chunk\_size})
fecharam aproximadamente $86\%$ do gap original que separava o gate da
Fase 2 ($0{,}515$) da meta de $0{,}80$.

O gap residual de $0{,}039$ tem três fatores complementares descritos na
literatura, ainda presentes no \textit{pipeline} mesmo após as intervenções.

\textbf{Comportamento abstrativo do gerador.} Modelos generativos como o
Llama~3.3~70B foram caracterizados desde \citeauthor{lewis2021}
\cite{lewis2021} como abstrativos: o decodificador sintetiza informação dos
\textit{chunks} em vez de reproduzi-los literalmente. Essa síntese é, conforme
o próprio paper seminal aponta, \textit{vantagem} para tarefas de
pergunta-resposta abertas, mas a métrica \textit{Faithfulness} do RAGAS
\cite{es2024} mede correspondência literal entre afirmações e contexto ---
penalizando o que pode ser paráfrase clinicamente correta. A ablação $2 \times 2$
(Seção~\ref{sec:ablacao-2x2}) já evidenciara esse efeito: substituir o Llama
por GPT-4o, com seu comportamento mais prompt-adherent, piorou
\textit{Faithfulness} em vez de melhorá-la.

\textbf{Granularidade do \textit{chunk}.} \citeauthor{gao2024} \cite{gao2024}
identifica o tamanho do \textit{chunk} como um dos principais determinantes
do \textit{trade-off} entre fidelidade literal e qualidade da resposta.
\citeauthor{wang2024bestpractices} \cite{wang2024bestpractices} sustenta
empiricamente que \textit{chunks} de 300--500 \textit{tokens}, combinados com
\textit{reranking}, otimizam esse trade-off em domínios técnicos. O teste
descrito na Seção~\ref{sec:chunk-test} reduziu o \texttt{CHUNK\_SIZE} de 800
para 500 \textit{chars} e não produziu ganho mensurável, porque o
\textit{chunker} hierárquico do \textit{pipeline} é dominado por fronteiras
de cabeçalho do \textit{Markdown} e pelo piso \texttt{MIN\_CHUNK\_SIZE = 400}.
Para ativar o efeito previsto pela literatura seria necessária a substituição
do \textit{chunker} por uma estratégia baseada em \textit{tokens} do modelo
(\textit{e.g.}, \textit{Small-to-Big} \cite{advancedChunkingRAG2025}), tarefa
registrada como trabalho futuro.

\textbf{Negação e nuances clínicas.} O domínio ILTB tem alta densidade de
regras condicionais (``tratar se X, exceto se Y'') e termos com semântica
oposta próxima (``elegível / contraindicado''), que \textit{retrievers}
densos atuais não capturam tão bem quanto modelos esparsos
\cite{cormack2009, formal2021}. A introdução de BM25 e do \textit{reranker}
mGTE (Seções~\ref{sec:hybrid} e \ref{sec:rerank}) mitigou parcialmente esse
limite, mas \textit{bi-encoders} dense permanecem suscetíveis a confundir
afirmações negadas e suas contrapartes afirmativas
\cite{huang2023halucination}.

A combinação desses três fatores residuais explica por que o \textit{pipeline}
estabiliza em $0{,}761$ em vez de atingir $0{,}80$. A próxima intervenção
prevista pela literatura --- \textit{self-reflection} sobre a resposta gerada
\cite{jeong2024selfbiorag} --- exige reformulação arquitetural com custo
desproporcional ao gap residual ($+8$ a $+15$\,pp esperados, ao preço de
dobrar as chamadas ao gerador), e está registrada como linha de pesquisa
prioritária no Capítulo~\ref{cap:conclusao}.

A revisão de \citeauthor{ocaf2024} \cite{ocaf2024} reforça ainda um ponto
metodológico relevante: a avaliação automatizada por RAGAS deve ser
\textbf{complementada} por avaliação humana estruturada, especialmente para
sistemas em que síntese clinicamente correta é desejável --- princípio
endossado também pela diretriz de relato \textit{TRIPOD-LLM}
\cite{gallifant2025tripod}. Esse argumento motiva o desenho da avaliação com
especialista descrita na Seção~\ref{sec:sessao-expert}, que complementa as
métricas automatizadas com uma rubrica clínica especializada e adiciona, em
particular, sensibilidade ao \textit{trade-off} segurança-vs-cobertura que
o RAGAS por si só não captura.

\section{Avaliação com Especialista}
\label{sec:sessao-expert}

A avaliação automatizada RAGAS, por suas limitações intrínsecas
(Seção~\ref{cap:referencial}), foi complementada por uma sessão estruturada com
especialista em ILTB. O delineamento dessa etapa segue as recomendações da literatura
de RAG biomédico, que reporta número médio de quatro avaliadores em estudos do
domínio --- variando de um a dez --- e enfatiza que avaliações com $N$ reduzido devem
ser estruturadas em torno de instrumentos multidimensionais, transformando o
especialista em consultor de validação técnica em lugar de sujeito estatístico
\cite{ocaf2024}.

\subsection{Perfil do Especialista}

\aluno{Caracterizar o expert: área de atuação (enfermagem em saúde pública, atenção
primária), tempo de experiência no manejo da ILTB, vínculo institucional (anonimizar
se necessário), formação. Incluir Termo de Consentimento Livre e Esclarecido
assinado.}

\subsection{Protocolo da Sessão}

A sessão foi estruturada em quatro etapas, alinhadas às atividades de
\textit{Demonstração} e \textit{Avaliação} do DSRM \cite{peffers2007}:

\begin{enumerate}[label=\textbf{(\arabic*)}]
  \item \textbf{Abertura e consentimento:} assinatura do TCLE, apresentação dos
        objetivos do chatbot, do escopo do corpus e dos limites declarados da
        ferramenta (apoio à decisão, não substituição do julgamento clínico).
  \item \textbf{Demonstração funcional:} apresentação das principais funcionalidades
        da interface, com exemplos previamente preparados, para que o especialista
        compreenda o fluxo de interação antes da fase avaliativa.
  \item \textbf{Interação com cenários clínicos:} o especialista submete consultas
        ao chatbot com base em perguntas reais que recebe na prática
        (rastreamento de PVHIV, esquemas terapêuticos, populações especiais,
        interações medicamentosas). Em paralelo, são apresentadas amostras das 38
        questões já avaliadas pelo RAGAS para revisão dirigida.
  \item \textbf{\textit{Debriefing} estruturado:} entrevista semi-estruturada
        orientada por instrumentos validados, descritos a seguir.
\end{enumerate}

Em conformidade com a política de privacidade descrita na
Seção~\ref{sec:lgpd-ref}, a sessão adota retenção zero do histórico de conversa,
\textit{logging} mínimo e ausência de dados identificáveis nas consultas.

\subsection{Instrumentos de Avaliação}

A avaliação combina quatro instrumentos, conforme recomendado para validação
multidimensional de chatbots clínicos com avaliador único \cite{ocaf2024}:

\begin{enumerate}[label=\textbf{I\arabic*}]
  \item \textbf{Acurácia técnica contra \textit{ground truth} do Ministério da Saúde:}
        o especialista classifica cada resposta como \textit{correta},
        \textit{parcialmente correta} ou \textit{incorreta} em relação ao protocolo
        oficial \cite{brasil2022}, indicando o trecho conflitante quando aplicável.
  \item \textbf{Classificação ordinal de severidade do erro:} escala Likert de cinco
        pontos (sem risco, risco baixo, risco moderado, risco alto, conduta
        perigosa) aplicada a cada resposta incorreta ou parcialmente correta,
        permitindo análise de segurança independente da acurácia bruta.
  \item \textbf{Tríade de qualidade RAG (rubrica qualitativa):} para cada resposta, o
        especialista avalia separadamente \textit{fidelidade} (a resposta deriva do
        contexto recuperado?), \textit{relevância da resposta} (responde à pergunta?)
        e \textit{relevância do contexto} (o trecho recuperado é o correto do
        protocolo?), seguindo a estrutura proposta por \citeauthor{es2024}
        \cite{es2024}.
  \item \textbf{TAM-AIN:} questionário fechado das quatro dimensões do
        \textit{Technology Acceptance Model for AI in Nursing}
        \cite{nursingTAM2024} --- alinhamento ético, prontidão organizacional,
        preservação da identidade profissional e capacidade de infraestrutura ---
        complementado por entrevista qualitativa sobre barreiras e facilitadores
        percebidos para adoção no SUS.
\end{enumerate}

A escolha por um instrumental multimétrico em detrimento de estatísticas tradicionais
de concordância inter-avaliador (e.g., coeficiente Kappa \cite{landis1977}) decorre
da natureza da amostra: com $N=1$, instrumentos como Kappa de Cohen perdem aderência
metodológica, pois exigem ao menos dois avaliadores para o cálculo da concordância
acima do acaso. A literatura defende, nesses cenários, o deslocamento do esforço
analítico para validade de conteúdo e segurança clínica \cite{ocaf2024}.

\subsection{Resultados e Análise da Sessão}

\aluno{Apresentar: número de questões avaliadas, distribuição de respostas por
categoria (correta / parcialmente correta / incorreta), severidade média dos erros,
scores TAM-AIN por dimensão, exemplos representativos de respostas avaliadas
positiva e negativamente, observações qualitativas do especialista.}

% ─────────────────────────────────────────────────────────────────────────────
\chapter{Discussão}
\label{cap:discussao}
% ─────────────────────────────────────────────────────────────────────────────

\section{Interpretação dos Resultados}

Os resultados consolidados da avaliação RAGAS, à luz da ablação $2 \times 2$
descrita na Seção~\ref{sec:ablacao-2x2} e da configuração final pós-migração
da Seção~\ref{sec:gate-pos-migracao}, comportam três leituras complementares.

\textbf{Primeira: o \textit{embedding} era o gargalo dominante.} A migração de
\texttt{paraphrase-multilingual-MiniLM-L12-v2} (384 dimensões) para
\texttt{BAAI/bge-m3} (1024 dimensões) produziu salto material em todas as
métricas para os dois geradores avaliados (Tabela~\ref{tab:ablacao-2x2}). Esse
achado é convergente com os \textit{benchmarks} multilíngues
\cite{SHTI2023, mmteb2025, marone2025mmbert} e com a revisão sistemática de
\citeauthor{ocaf2024} \cite{ocaf2024}, que identifica a escolha do
\textit{embedding} como o principal lever de qualidade em RAG clínico. O
patamar atingido de \textit{Context Precision} ($0{,}949$ na configuração
final) ultrapassa confortavelmente a meta de $0{,}75$, validando empiricamente
a hipótese predita pela literatura para o contexto deste trabalho.

\textbf{Segunda: modelo maior não dissolve o gap.} A substituição do
Llama~3.3~70B pelo GPT-4o, mantendo o restante do \textit{pipeline} constante,
\emph{piorou} a \textit{Faithfulness} em ambas as configurações de
\textit{embedding} ($0{,}515 \to 0{,}383$ no MiniLM; $0{,}675 \to 0{,}600$ no
BGE-M3). O GPT-4o adere mais estritamente às instruções do \textit{system
prompt} e produz respostas de \textit{fallback} (recusa) com maior frequência,
comportamento que o RAGAS pontua como menos fiel por ausência de \textit{claims}
verificáveis. Achados análogos em RAG clínico são descritos por
\citeauthor{zakka2024almanac} \cite{zakka2024almanac} no sistema Almanac
(NEJM AI). Este paradoxo tem duas implicações: (i) a escolha do Llama~3.3~70B
como gerador do artefato é tecnicamente justificada pelos próprios números, não
apenas por considerações de custo (Groq oferece o modelo a custo zero no
\textit{free tier}); (ii) métricas automatizadas como o RAGAS exigem
complementação por avaliação humana \cite{ocaf2024, gallifant2025tripod}, pois
podem penalizar comportamento desejável do ponto de vista de segurança
clínica.

\textbf{Terceira: intervenções pós-recuperação fecharam $\sim 86\%$ do gap
residual de \textit{Faithfulness}.} A ablação $2 \times 2$ identificou que o
\textit{retriever} era o gargalo dominante, mas a migração para BGE-M3 ainda
deixou \textit{Faithfulness} em $0{,}675$. Três iterações subsequentes,
fundamentadas na literatura específica e descritas nas
Seções~\ref{sec:hybrid}--\ref{sec:gate-final}, elevaram esse patamar para
$0{,}761$: (i) busca híbrida BM25 + denso com \textit{Reciprocal Rank Fusion}
\cite{cormack2009, formal2021, khattab2020}, contribuição $+0{,}014$ na
\textit{Faithfulness} e $+0{,}087$ na \textit{Context Recall};
(ii) \textit{reranking} \textit{cross-encoder} mGTE
\cite{zhang2024mgte, wang2024bestpractices}, contribuição $+0{,}072$ na
\textit{Faithfulness} e recuperação completa de \textit{Context Precision}
(para $0{,}964$); (iii) teste de \textit{chunk\_size} reduzido sem efeito
mensurável (Seção~\ref{sec:chunk-test}), limitado pela estrutura hierárquica
do \textit{chunker}. O gap residual final de $0{,}039$ converge ao patamar
descrito por \citeauthor{ocaf2024} \cite{ocaf2024} como teto típico antes de
intervenções de \textit{self-reflection} \cite{jeong2024selfbiorag}, registrada
como linha de pesquisa imediata no Capítulo~\ref{cap:conclusao}.

A leitura conjunta dos quatro \textit{scores} é coerente com o cenário previsto
pela revisão sistemática de \citeauthor{ocaf2024} \cite{ocaf2024}: sistemas
RAG biomédicos obtêm ganho médio de $1{,}35\times$ em eficácia clínica sobre
LLMs sem recuperação, mas atingem um teto comportamental quando a métrica
exige fidelidade verbatim a um modelo abstrativo. A resposta metodológica
adequada não é forçar o LLM a respostas extrativas (que diminuiria a utilidade
clínica em consultas que exigem síntese), mas combinar duas estratégias
complementares: (i) avançar tecnicamente via intervenções pós-recuperação
referenciadas acima; (ii) complementar a avaliação com instrumento humano
estruturado, conforme detalhado na Seção~\ref{sec:sessao-expert}.

\section{Limitações}

\begin{enumerate}[label=\textbf{L\arabic*}]
  \item Corpus limitado a documentos públicos --- nenhum dado clínico real de
        pacientes foi utilizado;
  \item Ausência de \textit{dataset} clínico anotado em PT-BR para validação
        RAGAS independente;
  \item Avaliação qualitativa com amostra reduzida (1 expert);
  \item Dependência de API externa para o LLM (Groq) --- a operação do
        \textit{free tier} (limite de 100K tokens diários, TPD) restringe
        ciclos experimentais; o presente trabalho recorreu ao plano
        \textit{Dev Tier} pago para concluir a célula Llama+BGE-M3 da
        ablação $2 \times 2$ (Seção~\ref{sec:ablacao-2x2});
  \item \textit{Pipeline} atual incorpora busca híbrida BM25+denso e
        \textit{reranking} \textit{cross-encoder} mGTE (Seções
        \ref{sec:hybrid} e \ref{sec:rerank}), mas não inclui
        \textit{self-reflection} sobre a resposta gerada
        \cite{jeong2024selfbiorag}. Essa intervenção, com expectativa
        documentada de $+8$ a $+15$\,pp em \textit{Faithfulness}, é a próxima
        linha de pesquisa para fechar o gap residual de $0{,}039$ até a meta
        de $0{,}80$ (Capítulo~\ref{cap:conclusao}). \textit{Chunking} baseado
        em \textit{tokens} (\textit{e.g.}, \textit{Small-to-Big}) é
        intervenção complementar potencial não testada
        \cite{advancedChunkingRAG2025}.
\end{enumerate}

\section{Conformidade com a LGPD}

O sistema não processa dados pessoais de pacientes: o corpus é composto
exclusivamente de documentos públicos e as interações com o chatbot não são
armazenadas. O modelo de embedding é executado localmente, sem envio de
\textit{queries} a APIs externas antes da geração \cite{lgpdSaude2023,
sokPrivacyLLM2026}.

A configuração adotada para esta prova de conceito é, contudo, uma escolha
deliberadamente conservadora. Em cenários de implantação institucional --- ainda fora
do escopo deste trabalho, mas previstos como continuidade natural ---, três pontos
de atenção emergem da literatura.

O primeiro é a \textbf{entrada de dados identificáveis pelo usuário final}. Mesmo com
corpus público e API sem armazenamento, o conteúdo das consultas formuladas pelo
profissional pode trazer informações clinicamente identificáveis (descrição de caso,
data, comorbidades, identificadores indiretos). \citeauthor{privacyRAGHealthcare2025}
\cite{privacyRAGHealthcare2025} categorizam o pesquisador e o profissional clínico
como \textit{adversários internos acidentais}: agentes legítimos cujo fluxo de
trabalho pode causar vazamento inadvertido. Para mitigar esse risco, recomenda-se
implementar \textit{scrubbing} ex-ante no \textit{frontend} --- redação automática de
identificadores antes do envio ao retriever --- e exibir aviso explícito ao usuário
sobre a política de tratamento de dados sensíveis.

O segundo é a \textbf{rastreabilidade auditável}. Sistemas de IA em saúde sob a LGPD
demandam \textit{logging} compatível com auditoria sem violar a finalidade declarada
do tratamento \cite{lgpdSaude2023, sokPrivacyLLM2026}. Em produção, isso implica
registrar metadados de uso (timestamps, tipo de consulta, presença de
\textit{disclaimer}, identificador de sessão pseudonimizado) sem reter o conteúdo das
consultas, mecanismo descrito como \textit{logging consciente de privacidade}
\cite{privacyRAGHealthcare2025}. A retenção zero adotada nesta POC não escala para um
serviço institucional, que precisa demonstrar rastreabilidade frente a eventuais
incidentes.

O terceiro é a \textbf{relação jurídica com o provedor do LLM}. O modelo Llama~3.3
70B foi acessado neste trabalho via API Groq, sem armazenamento declarado de prompts;
em ambiente institucional, contudo, recomenda-se a celebração de Acordo de
Tratamento de Dados (DPA, na sigla em inglês) com o operador, em conformidade com os
artigos 39 e 47 da LGPD, ou o emprego de \textit{inference} local em GPU dedicada
para eliminar a transferência. A literatura específica sobre RAG biomédico endossa
essa estratégia, observando que adversários do lado da nuvem podem reconstruir
trechos sensíveis a partir de ativações intermediárias mesmo quando o provedor não
armazena prompts explicitamente \cite{privacyRAGHealthcare2025, sokPrivacyLLM2026}.

Em síntese, o desenho atual atende à LGPD por exclusão de dados sensíveis, e não por
mecanismos ativos de proteção. Esses mecanismos são pré-requisito de qualquer
implantação real e estão contemplados como linha de pesquisa subsequente
(Capítulo~\ref{cap:conclusao}).

\section{Implicações para Adoção}

A adoção efetiva de chatbots clínicos como o desenvolvido neste trabalho depende de
fatores que ultrapassam a qualidade técnica do artefato. O modelo TAM-AIN
\cite{nursingTAM2024}, adotado como referência para a avaliação com especialista
(Seção~\ref{sec:sessao-expert}), aponta quatro dimensões críticas que devem ser
ponderadas em qualquer estratégia de implantação no SUS.

\textbf{Alinhamento ético e confiança.} A literatura indica padrão de
\textit{otimismo cauteloso} entre profissionais de enfermagem frente à IA: cerca de
85\% reconhecem benefícios potenciais para precisão diagnóstica e redução de erro
humano, mas aproximadamente 70\% mantêm reservas quanto ao componente humano do
cuidado e à segurança dos dados \cite{nursingTAM2024}. Para o chatbot ILTB, isso
implica investir em comunicação clara dos limites da ferramenta e da política de
privacidade --- sob pena de minar a base de confiança necessária para uso continuado.

\textbf{Preservação da identidade profissional.} O TAM-AIN distingue ``ferramenta de
apoio'' de ``substituto da decisão clínica'', dimensão particularmente sensível
entre profissionais com mais tempo de prática, que tendem a manifestar receio de
desqualificação. A interface da POC já incorpora \textit{disclaimer} explícito de
não substituição do julgamento profissional, mas a literatura sugere que a aderência
de longo prazo exige que o sistema demonstre, no fluxo da interação, capacidade de
reforçar o protagonismo do enfermeiro --- por exemplo, citando fonte e contexto a
cada resposta para que o profissional valide ativamente as informações.

\textbf{Prontidão organizacional e infraestrutura.} A heterogeneidade do parque
tecnológico das unidades de atenção primária do SUS configura barreira
infraestrutural concreta: conexões intermitentes, navegadores desatualizados, falta
de dispositivos dedicados \cite{nursingTAM2024}. A arquitetura web atual, leve e
sem dependência de instalação local no \textit{client}, é compatível com esse cenário,
mas estudos futuros precisarão verificar latência percebida em redes de menor
qualidade e estratégias de \textit{caching} para uso \textit{offline} parcial.

\textbf{Carga de trabalho e usabilidade percebida.} Entre as barreiras mais
reportadas pela literatura estão a sobrecarga e o ceticismo sobre o tempo
adicional exigido para incorporar novas ferramentas ao fluxo. Avaliações
subsequentes deverão medir, além da acurácia clínica das respostas, o ganho líquido
de tempo na consulta a protocolos --- métrica operacional que sustenta o argumento
de utilidade frente ao gestor da unidade.

Esses fatores indicam que a transição da POC para implantação institucional não
constitui apenas tarefa de engenharia (resolvendo as questões discutidas na
Seção~\ref{cap:discussao}), mas também tarefa organizacional, dependente de adesão
da liderança local e de capacitação dos profissionais.

% ─────────────────────────────────────────────────────────────────────────────
\chapter{Conclusão}
\label{cap:conclusao}
% ─────────────────────────────────────────────────────────────────────────────

Este trabalho desenvolveu e avaliou um chatbot de apoio clínico ao manejo da
Infecção Latente por \textit{Mycobacterium tuberculosis}, implementado com a
arquitetura RAG e documentos oficiais em língua portuguesa, seguindo a
metodologia DSRM \cite{peffers2007}. O artefato é composto por um \textit{pipeline}
local de \textit{embeddings} multilíngues, um \textit{vector store} ChromaDB com
indexação HNSW e o modelo Llama~3.3~70B acessado via API Groq como gerador, e
foi disponibilizado por meio de uma interface web servida pelo FastAPI.

A avaliação automatizada com o \textit{framework} RAGAS sobre 38 questões
clínicas foi conduzida em três etapas. A configuração inicial, com
\textit{embedding} 384D e o Llama~3.3~70B como gerador, atingiu
\textit{Context Precision} de $0{,}735$ e \textit{Faithfulness} de $0{,}515$
--- patamar consistente com a literatura de RAG biomédico, mas insuficiente
para a meta de $0{,}80$. Uma ablação $2 \times 2$ cruzando gerador (Llama vs
GPT-4o) e \textit{embedding} (MiniLM 384D vs BGE-M3 1024D) isolou o
\textit{embedding} como gargalo dominante; a migração para
\texttt{BAAI/bge-m3} elevou todas as métricas a um patamar intermediário com
\textit{Faithfulness} de $0{,}675$. Três iterações pós-recuperação
subsequentes --- busca híbrida BM25 + denso com \textit{Reciprocal Rank
Fusion}, \textit{reranking} \textit{cross-encoder} mGTE e teste de
\textit{chunk\_size} reduzido --- consolidaram o gate operacional final do
artefato em \textit{Faithfulness} $0{,}761$ (a $0{,}039$ da meta original),
\textit{Context Precision} $0{,}964$ (PASS robusto), \textit{Answer
Relevancy} $0{,}720$ e \textit{Context Recall} $0{,}803$, um salto cumulativo
de $+47{,}8\%$ na \textit{Faithfulness} desde o \textit{baseline} da Fase 2.
Adicionalmente, a substituição do Llama pelo GPT-4o piorou a
\textit{Faithfulness} em ambos os \textit{embeddings} ($0{,}515 \to 0{,}383$
no MiniLM; $0{,}675 \to 0{,}600$ no BGE-M3), achado convergente com
observações independentes em RAG clínico \cite{zakka2024almanac} e que
justifica tecnicamente a escolha do Llama como gerador do artefato.
\aluno{Completar com 2--3 frases sobre o resultado da sessão estruturada
com o especialista (acurácia clínica média, \textit{scores} TAM-AIN,
principais observações qualitativas) após a execução da
Seção~\ref{sec:sessao-expert}.}

As principais contribuições são o artefato aberto, o \textit{dataset} de 38
questões clínicas em português sobre ILTB com \textit{ground truths} extraídos
verbatim dos protocolos, a análise empírica do \textit{trade-off} LGPD vs.
desempenho e a proposta de protocolo de avaliação multidimensional com $N=1$
combinando RAGAS, rubrica clínica e TAM-AIN. \aluno{Ao finalizar, ajustar este
parágrafo se a sessão com expert produzir contribuição empírica adicional.}

\section{Trabalhos Futuros}
\label{sec:trabalhos-futuros}

A avaliação experimental e a revisão bibliográfica conduzidas no escopo deste
trabalho identificaram um conjunto de direções de evolução, agrupadas em
três frentes. A frente técnica foi parcialmente endereçada no próprio
desenvolvimento do artefato (Seções \ref{sec:hybrid} e \ref{sec:rerank}, com
ganho cumulativo de $+0{,}086$ em \textit{Faithfulness}); restam linhas de
pesquisa para fechar o gap residual final
($0{,}761 \to \geq 0{,}80$).

\textbf{Frente técnica:}

\begin{enumerate}[label=\textbf{TF\arabic*}]
  \item \textbf{Self-RAG / \textit{self-reflection}} --- após gerar a
        resposta, um \textit{critic step} avalia se está suportada pelos
        \textit{chunks} e dispara nova recuperação ou refinamento quando
        necessário. \citeauthor{jeong2024selfbiorag} \cite{jeong2024selfbiorag}
        reportam ganho de $+8$ a $+15$ pontos em \textit{Faithfulness} em
        QA biomédico \textit{long-form}, ao custo de aproximadamente dobrar
        o número de chamadas ao LLM por consulta. Esta é a intervenção mais
        promissora para fechar o gap residual final de $0{,}039$ até a meta
        de $0{,}80$;
  \item \textbf{\textit{Chunking} estrito por \textit{tokens} do modelo} ---
        substituir o \textit{chunker} hierárquico atual (guiado por
        cabeçalhos do \textit{Markdown}) por estratégia tipo
        \textit{Small-to-Big}, \textit{Sliding Window} ou proposições
        atômicas (DenseX) \cite{advancedChunkingRAG2025, gao2024},
        permitindo o efeito de \textit{chunks} de 300--500 \textit{tokens}
        previsto por \citeauthor{wang2024bestpractices}
        \cite{wang2024bestpractices} mas não capturado no teste descrito na
        Seção~\ref{sec:chunk-test};
  \item \textbf{Prompt v5 com citação granular de \textit{span}} ---
        instruir o gerador a citar não apenas o documento, mas o trecho
        literal de cada \textit{claim}. \citeauthor{zakka2024almanac}
        \cite{zakka2024almanac} usa essa estratégia no sistema Almanac;
        ganho esperado em fidelidade, ao risco de intensificar
        comportamento conservador do gerador (a história de prompts v2--v4
        descrita no Apêndice~\ref{ape:prompt} mostra que prompts mais
        restritivos podem piorar a métrica RAGAS).
\end{enumerate}

\textit{Já incorporadas no escopo deste trabalho:} migração de
\textit{embedding} para BGE-M3 1024D (Seção~\ref{sec:gate-pos-migracao}),
busca híbrida BM25 + denso com RRF (Seção~\ref{sec:hybrid}) e
\textit{reranking} \textit{cross-encoder} mGTE (Seção~\ref{sec:rerank}).

\textbf{Frente metodológica:}

\begin{enumerate}[label=\textbf{TF\arabic*}, start=4]
  \item \textbf{Piloto clínico com múltiplos profissionais de enfermagem} ---
        com aprovação de Comitê de Ética em Pesquisa (CEP), permitindo
        cálculo de concordância inter-avaliador (e.g., Kappa
        \cite{landis1977}) e validação populacional do instrumento TAM-AIN
        \cite{nursingTAM2024};
  \item \textbf{Construção de \textit{benchmark} público em português} ---
        consolidação do \textit{dataset} de 38 questões com revisão por painel
        de especialistas, em formato compatível com avaliações futuras de
        outros sistemas RAG clínicos em PT-BR.
\end{enumerate}

\textbf{Frente operacional:}

\begin{enumerate}[label=\textbf{TF\arabic*}, start=6]
  \item \textbf{Implantação institucional} --- servidor dedicado com HTTPS,
        autenticação, \textit{logging} consciente de privacidade e Acordo de
        Tratamento de Dados (DPA) com o operador do LLM, conforme discutido na
        Seção~\ref{cap:discussao};
  \item \textbf{Canais de acesso ampliados} --- integração com sistemas
        existentes do SUS (e.g., e-SUS APS) e canal via WhatsApp Business API,
        aderente à expectativa de uso real por enfermeiros em campo;
  \item \textbf{Inferência local em \textit{hardware} dedicado} --- migração do
        gerador para servidor com GPU própria, eliminando a transferência de
        \textit{prompts} a operador externo e fechando integralmente o
        \textit{loop} de conformidade LGPD descrito no
        Capítulo~\ref{cap:discussao}.
\end{enumerate}

% ═════════════════════════════════════════════════════════════════════════════
% ELEMENTOS PÓS-TEXTUAIS
% ═════════════════════════════════════════════════════════════════════════════
\postextual

\bibliography{referencias}

% ─── Apêndice A — System Prompt ──────────────────────────────────────────────
\begin{apendicesenv}

\partapendices

\chapter{System Prompt de Produção (v1)}
\label{ape:prompt}

A versão~v1 do \textit{system prompt} foi a que produziu o melhor resultado de
\textit{Faithfulness} entre as quatro testadas
(v1~$=0{,}586$; v4~$=0{,}574$; v3~$=0{,}457$; v2~$=0{,}429$), e foi mantida como
configuração de produção. As versões alternativas v2 e v3 introduziam restrições
fortes (``EXCLUSIVAMENTE'', limite de quatro frases, citação obrigatória de
documento e seção) que paradoxalmente reduziram a métrica RAGAS, pois levavam o
modelo a fazer \textit{fallbacks} desnecessários ou a incluir referências
explícitas a fontes que não constavam dos \textit{chunks} recuperados. A versão
v4, baseada em \textit{few-shot} com dois exemplos autocontidos, ficou ligeiramente
abaixo de v1. A análise detalhada dessas iterações consta no diário técnico do
projeto.

\begin{quote}
\begin{small}
\begin{verbatim}
Você é um assistente clínico especializado nos protocolos do
Ministério da Saúde para Infecção Latente pelo Mycobacterium
tuberculosis (ILTB).

Regras obrigatórias:
1. Responda SOMENTE com base nos trechos de protocolo fornecidos
   no contexto.
2. Se a informação não estiver no contexto, diga: "Não encontrei
   essa informação nos protocolos indexados. Consulte o Manual
   de Recomendações do MS."
3. Cite sempre a seção de origem da resposta.
4. Não faça diagnósticos nem prescrições — apenas forneça
   informação de protocolo.
5. Use linguagem técnica adequada para enfermeiros.
\end{verbatim}
\end{small}
\end{quote}

O contexto recuperado pelo \textit{retriever} é injetado pelo \textit{user
message}, não pelo \textit{system prompt}, conforme implementação em
\texttt{app/src/llm/client.py}. Os \textit{chunks} são apresentados ao modelo na
ordem decrescente de similaridade cossenoidal, com identificador de fonte
(documento e cadeia de cabeçalhos) anexado a cada um.
de demonstração.}

% ─── Apêndice B — Dataset de Avaliação ───────────────────────────────────────
\chapter{Dataset de Avaliação RAGAS}
\label{ape:dataset}

O \textit{dataset} de avaliação contém 42 questões (38 \textit{in-scope} + 4
\textit{out-of-scope}), com \textit{ground truths} extraídos verbatim dos
arquivos Markdown indexados a partir dos protocolos do MS e da OMS. A tabela
abaixo lista todas as questões, agrupadas por categoria
(ET, IND, PE, DI, MO, IT, EA, FE).

{\footnotesize
\begin{longtable}{@{}p{1.4cm}@{\hspace{1mm}}p{0.7cm}@{\hspace{1mm}}p{4.5cm}p{6.5cm}@{}}
\toprule
\textbf{ID} & \textbf{Cat.} & \textbf{Pergunta} & \textbf{\textit{Ground truth}} \\
\midrule
\endfirsthead
\multicolumn{4}{c}{\textit{(continuação da Tabela do Apêndice~\ref{ape:dataset})}} \\
\toprule
\textbf{ID} & \textbf{Cat.} & \textbf{Pergunta} & \textbf{\textit{Ground truth}} \\
\midrule
\endhead
\midrule
\multicolumn{4}{r}{\textit{Continua na próxima página}} \\
\endfoot
\bottomrule
\endlastfoot
\texttt{ET-01} & ET & Qual a dose de isoniazida no esquema 3HP para adultos? & No esquema 3HP, adultos (>14 anos, ≥30kg) tomam 900mg de isoniazida por semana, uma vez por semana, durante 3 meses (12 doses). O esquema prevê ajuste de dose por peso corporal. \\ \hline
\texttt{ET-02} & ET & Qual a duração do esquema 3HP e quantas doses são necessárias? & O tratamento 3HP estará completo quando ocorrer a tomada de 12 doses de isoniazida + rifapentina por 12 semanas. Dependendo do caso, esse prazo pode ser prorrogado para 15 semanas. \\ \hline
\texttt{ET-03} & ET & Qual a dose de rifampicina no esquema 4R para adultos e por quanto tempo? & Rifampicina 600 mg/dia, indicada para adultos ≥ 50 kg, administrada diariamente por 4 meses, totalizando 120 doses; que podem ser completadas no período de 4 a 6 meses, conforme adesão ao tratamento. \\ \hline
\texttt{ET-04} & ET & Qual a dose de isoniazida no esquema 6H para crianças? & Esquemas pediátricos incluem isoniazida (10 mg/kg/dia; dose máxima: 300 mg/dia) administrada diariamente por 6 meses (totalizando 180 doses, a serem completadas em 6 a 9 meses). \\ \hline
\texttt{ET-05} & ET & Qual esquema é preferencial para hepatopatas e idosos acima de 50 anos? & A rifampicina é o tratamento de escolha para hepatopatas e pessoas acima de 50 anos. \\ \hline
\texttt{ET-06} & ET & Qual a dose pediátrica de isoniazida no esquema 3HP para crianças de 10 a 15 kg? & No esquema 3HP pediátrico (crianças de 2 a 14 anos), a dose de isoniazida para crianças de 10 a 15 kg é de 300mg por semana. Para rifapentina, a dose é também 300mg por semana para esse peso. \\ \hline
\texttt{ET-07} & ET & Qual a dose de isoniazida para adultos no esquema 6H? & Para adultos, a dose de isoniazida nos esquemas 6H e 9H é de 5 a 10 mg/kg/dia, com dose máxima de 300 mg/dia. O esquema 6H contempla 180 doses (6 a 9 meses) e o 9H contempla 270 doses (9 a 12 meses). O esquema 9H protege mais que o 6H e deve ser preferido sempre que houver boa adesão. \\ \hline
\texttt{IT-01} & IND & Quem deve receber tratamento da ILTB independente do resultado da PT ou IGRA? & Alguns grupos devem ser tratados sem PT e sem IGRA, como contatos de alto risco com TB bacilífera. O Quadro 1 do Protocolo de Vigilância da ILTB especifica as indicações de tratamento, incluindo grupos que devem tratar sem PT/IGRA, grupos que tratam com PT ≥ 5mm ou IGRA positivo, e grupos que tratam com PT ≥ 10mm ou IGRA positivo. \\ \hline
\texttt{IT-02} & IND & Quando é indicado tratar ILTB em pessoa com alteração radiológica fibrótica pulmonar? & Todos os indivíduos com PPD maior ou igual a 5mm e IGRA positivo e que tenham alterações radiológicas pulmonares fibróticas sugestivas de sequelas de TB, independe de ter alguma doença ou estar em tratamento, também têm indicação de tratamento da TBi. \\ \hline
\texttt{IT-03} & IND & Qual a indicação de tratamento da ILTB em paciente em pré-transplante? & Estão incluídos na indicação de tratamento da TBi os pacientes com PPD maior ou igual a 5mm ou IGRA positivo em pré-transplante em uso de terapia imunossupressora. \\ \hline
\texttt{IT-04} & IND & Quando tratar ILTB em paciente com prednisona? & Devem receber tratamento para TBi indivíduos com PPD maior ou igual a 5mm e IGRA positivo, que estejam em uso ou iniciarão terapia com dose acima de 15 mg/dia de prednisona por período superior a um mês. \\ \hline
\texttt{IT-05} & IND & Quando é indicado tratar ILTB em pessoas em contato com TB? & Recomenda-se o tratamento da infecção latente pelo M. tuberculosis (ILTB) para pessoas infectadas, identificadas por meio da prova tuberculínica (PT) ou Interferon-Gamma Release Assays (IGRA), quando apresentam risco de desenvolver TB. Antes do tratamento deve-se afastar definitivamente a tuberculose ativa. \\ \hline
\texttt{PE-01} & PE & Posso tratar ILTB em gestante HIV negativa? & Em gestantes HIV negativas, a recomendação é postergar o tratamento da ILTB para após o parto. \\ \hline
\texttt{PE-02} & PE & Quando tratar ILTB em gestante HIV positiva? & Em gestantes HIV positivas, preconiza-se o tratamento após o 3º mês de gestação. \\ \hline
\texttt{PE-03} & PE & Como manejar paciente em uso de anti-TNF com diagnóstico de ILTB? & Indivíduos com PPD ≥ 5mm e IGRA positivo que estejam em uso ou iniciarão terapia biológica com ação anti-TNF alfa têm indicação de tratamento para TBi. O tratamento deve ser iniciado 30 dias antes do uso da terapia imunossupressora. Após o diagnóstico durante o tratamento da DII, recomenda-se tratamento da TBi além de vigilância e monitoramento. \\ \hline
\texttt{PE-04} & PE & Qual esquema de tratamento da ILTB deve ser evitado em gestantes? & O esquema 3HP (Rifapentina + Isoniazida) não está indicado para gestantes. As indicações do 3HP excluem gestantes entre as populações para as quais o esquema pode ser utilizado. \\ \hline
\texttt{PE-05} & PE & Qual a dose de rifampicina para crianças no esquema 4R? & No esquema pediátrico 4R, a rifampicina é administrada na dose de 15 mg/kg/dia, com dose máxima de 600 mg/dia, administrada diariamente por 4 meses. \\ \hline
\texttt{PE-06} & PE & Qual o risco aumentado de TB doença em pacientes com DII em uso de anti-TNF? & Estudo realizado em nosso meio demonstrou aumento significativo do risco de TB doença em pacientes com DII em tratamento. O uso de tiopurinas aumentou o risco em 5,85 vezes, anti-TNF em 3,93 vezes e a combinação azatioprina + anti-TNF em 9,03 vezes. \\ \hline
\texttt{PE-07} & PE & Qual esquema usar para paciente com HIV? & Para PVHIV, o esquema preferencial é a isoniazida (9H), pois o esquema 4R (rifampicina) está contraindicado em PVHIV em uso de inibidores de protease e integrase. O esquema 9H (270 doses, 9 a 12 meses) protege mais que o 6H. PVHIV com contagem de CD4 ≤ 350 células/µL têm indicação de tratamento para ILTB independentemente do resultado da PT ou IGRA. \\ \hline
\texttt{DI-01} & DI & Qual o ponto de corte da prova tuberculínica para indicar tratamento da ILTB em pacientes em uso de anti-TNF? & Devem receber tratamento para TBi indivíduos com PPD maior ou igual a 5 mm e IGRA positivo, que estejam em uso ou iniciarão terapia biológica com ação anti-TNF alfa ou uso de dose acima de 15 mg/dia de prednisona por período superior a um mês. \\ \hline
\texttt{DI-02} & DI & Qual exame é superior para diagnóstico da ILTB: IGRA ou PPD? & Não há superioridade do IGRA sobre o PPD e vice-versa. O PPD sofre influência da vacina BCG até dois anos depois de sua aplicação. No caso do IGRA, existe a possibilidade de resultado indeterminado, especialmente em condições de moradia em região endêmica, tratamento com corticoide ou imunossupressor, atividade de doença, idade avançada, baixos níveis de albumina, linfopenia e altos níveis de proteína C reativa. \\ \hline
\texttt{DI-03} & DI & O que deve ser feito antes de iniciar o tratamento da ILTB? & Antes de efetuar o tratamento da ILTB deve-se: (a) afastar definitivamente a tuberculose ativa e (b) realizar a notificação por meio da Ficha de notificação das pessoas em tratamento da ILTB. O diagnóstico de ILTB deve ser realizado por meio dos testes da PT ou do IGRA. \\ \hline
\texttt{DI-04} & DI & Quando o resultado do IGRA for indeterminado, o que fazer? & Quando o resultado do IGRA for indeterminado, o exame deve ser repetido o mais breve possível. Se o resultado for mantido, é recomendado considerar o tratamento de TBi. \\ \hline
\texttt{DI-05} & DI & Qual a conduta para excluir TB ativa antes de tratar ILTB? & Para excluir TB ativa antes de tratar ILTB, deve-se realizar: anamnese, exame clínico, laboratorial e radiografia de tórax (por vezes tomografia computadorizada de tórax). É importante excluir todas as formas clínicas da doença (pulmonar/laríngea e extrapulmonar), além de considerar vínculos epidemiológicos e exames bacteriológicos. \\ \hline
\texttt{MO-01} & MO & Com que frequência devem ser realizadas as avaliações clínicas no esquema 3HP? & No esquema com isoniazida e rifapentina (3HP) estão indicadas avaliações mensais para avaliação clínica, seguimento da adesão ao tratamento e de eventos adversos. \\ \hline
\texttt{MO-02} & MO & Quando devo suspender o tratamento da ILTB por hepatotoxicidade? & Durante o tratamento da ILTB, a hepatite aguda medicamentosa (CID K71) é uma das causas de interesse para investigação e registro. O tratamento deve ser suspenso diante de sinais de hepatotoxicidade clinicamente significativa, como elevação de transaminases com sintomas. \\ \hline
\texttt{MO-03} & MO & O que considerar durante o acompanhamento do tratamento da ILTB além do tempo? & Durante o tratamento de TBi é importante levar em consideração o número de doses utilizadas e não apenas o tempo de tratamento. \\ \hline
\texttt{MO-04} & MO & Quais informações devem ser registradas ao final de cada consulta de acompanhamento? & Ao final da consulta de seguimento, deve-se: agendar consulta de seguimento para monitorar o progresso do tratamento, providenciar uma forma de o paciente entrar em contato com um profissional de saúde caso surjam problemas, organizar consultas em grupo se houver, registrar o que aconteceu durante a consulta, e encaminhar para serviços sociais existentes. \\ \hline
\texttt{MO-05} & MO & Com que frequência pacientes com DII em uso de imunossupressor devem ser rastreados para ILTB? & Pacientes com DII em uso contínuo de terapia com ação imunossupressora devem realizar rastreio anual para TBi por meio do PPD e/ou IGRA. Além disso, o consenso recomendou que pacientes sem tratamento prévio de TB realizem exames para triagem de TBi anuais durante o período de 3 anos e, em caso de mudança da terapia, a triagem está recomendada por um período de mais 3 anos. \\ \hline
\texttt{IM-01} & IT & Rifampicina tem interação com contraceptivos orais? & A rifampicina interage com contraceptivos orais, reduzindo sua eficácia por indução enzimática. É recomendado utilizar método alternativo de contracepção durante o uso da rifampicina. \\ \hline
\texttt{IM-02} & IT & Qual a principal preocupação ao usar rifampicina em pessoas vivendo com HIV em uso de antirretrovirais? & A rifampicina é um potente indutor enzimático que pode reduzir significativamente os níveis plasmáticos de antirretrovirais, especialmente inibidores de protease e inibidores não nucleosídeos da transcriptase reversa. Deve-se verificar interações medicamentosas antes de iniciar rifampicina em PVHIV. \\ \hline
\texttt{IM-03} & IT & Isoniazida tem interação com fenitoína? & A isoniazida interage com fenilhidantoína (fenitoína), causando maior hepatotoxicidade. A recomendação do Ministério da Saúde é evitar o uso concomitante. Quando ambos os fármacos forem necessários, deve-se monitorar sintomas e enzimas hepáticas conforme indicado. \\ \hline
\texttt{IM-04} & IT & Quanto tempo antes do início de terapia imunossupressora deve ser iniciado o tratamento da ILTB? & O Ministério da Saúde recomenda que o tratamento da TBi seja iniciado 30 dias antes do uso da terapia imunossupressora. \\ \hline
\texttt{IM-05} & IT & Qual o prazo recomendado para aguardar antes de iniciar tratamento imunossupressor após diagnóstico de ILTB, segundo o consenso europeu? & O Consenso Europeu que aborda infecções em DII recomenda aguardar 4 semanas para início do tratamento imunossupressor após o diagnóstico de ILTB, exceto em casos de maior urgência clínica após aconselhamento especializado. \\ \hline
\texttt{EA-01} & EA & Quais são os principais efeitos adversos da isoniazida? & Os principais efeitos adversos da isoniazida incluem: hepatotoxicidade (hepatite aguda medicamentosa), neuropatia periférica (pode ser prevenida com piridoxina), reações de hipersensibilidade. A hepatotoxicidade é a reação adversa mais grave e requer monitoramento das transaminases. \\ \hline
\texttt{EA-02} & EA & Como prevenir a neuropatia periférica causada pela isoniazida? & A neuropatia periférica causada pela isoniazida pode ser prevenida com a suplementação de piridoxina (vitamina B6). Em gestantes HIV negativas com indicação de tratamento, a isoniazida deve ser acompanhada de suplementação de piridoxina. \\ \hline
\texttt{EA-03} & EA & Quais doenças hepáticas são causas de encerramento do tratamento por efeito adverso? & As causas básicas potencialmente relacionadas à hepatotoxicidade de interesse para investigação e registro no IL-TB incluem: Tuberculose (CID A15 a A19), Complicações do HIV (CID B22, B22.7, B23, B23.8), Hepatite aguda medicamentosa (CID K71) e Doença hepática aguda viral (CID B17). \\ \hline
\texttt{EA-04} & EA & O que acontece com o tratamento da ILTB quando a paciente engravida durante o tratamento? & O diagnóstico de gravidez no decorrer do tratamento é uma causa de descontinuidade do tratamento. A paciente deverá iniciar um novo tratamento para ILTB após o parto. Para o novo tratamento, cada caso necessitará ser avaliado individualmente, considerando o número de doses tomadas até a interrupção do tratamento. \\ \hline
\texttt{FE-01} & FE & Qual o tratamento para tuberculose ativa com resistência à rifampicina? & --- \\ \hline
\texttt{FE-02} & FE & Qual a dose de amoxicilina para tratar pneumonia bacteriana? & --- \\ \hline
\texttt{FE-03} & FE & Quais são os critérios diagnósticos para COVID-19? & --- \\ \hline
\texttt{FE-04} & FE & Como fazer o esquema RIPE para tuberculose ativa? & --- \\
\end{longtable}
}

\end{apendicesenv}

\end{document}
